%File: formatting-instructions-latex-2023.tex
%release 2023.0
\documentclass[letterpaper]{article} % DO NOT CHANGE THIS
\usepackage{aaai23}  % DO NOT CHANGE THIS
\usepackage{times}  % DO NOT CHANGE THIS
\usepackage{helvet}  % DO NOT CHANGE THIS
\usepackage{courier}  % DO NOT CHANGE THIS
\usepackage[hyphens]{url}  % DO NOT CHANGE THIS
\usepackage{graphicx} % DO NOT CHANGE THIS
\urlstyle{rm} % DO NOT CHANGE THIS
\def\UrlFont{\rm}  % DO NOT CHANGE THIS
\usepackage{natbib}  % DO NOT CHANGE THIS AND DO NOT ADD ANY OPTIONS TO IT
\usepackage{caption} % DO NOT CHANGE THIS AND DO NOT ADD ANY OPTIONS TO IT
\usepackage{amssymb}
\usepackage{amsmath}
\usepackage{mathrsfs}
\frenchspacing  % DO NOT CHANGE THIS
\setlength{\pdfpagewidth}{8.5in}  % DO NOT CHANGE THIS
\setlength{\pdfpageheight}{11in}  % DO NOT CHANGE THIS
%
% These are recommended to typeset algorithms but not required. See the subsubsection on algorithms. Remove them if you don't have algorithms in your paper.
\usepackage{algorithm}
\usepackage{algorithmic}
\renewcommand{\phi}{\varphi}
\newcommand{\proves}{\vdash}
%
% These are are recommended to typeset listings but not required. See the subsubsection on listing. Remove this block if you don't have listings in your paper.
\usepackage{newfloat}
\usepackage{listings}
\DeclareCaptionStyle{ruled}{labelfont=normalfont,labelsep=colon,strut=off} % DO NOT CHANGE THIS
\lstset{%
	basicstyle={\footnotesize\ttfamily},% footnotesize acceptable for monospace
	numbers=left,numberstyle=\footnotesize,xleftmargin=2em,% show line numbers, remove this entire line if you don't want the numbers.
	aboveskip=0pt,belowskip=0pt,%
	showstringspaces=false,tabsize=2,breaklines=true}
\floatstyle{ruled}
\newfloat{listing}{tb}{lst}{}
\floatname{listing}{Listing}
\pdfinfo{
/TemplateVersion (2023.1)
}

\setcounter{secnumdepth}{0}
\title{Progress Report --- Agora Forge: Synthesizing Multiagent Mathematical Markets}
\author{
    Riyaz Ahuja,
    Shivansh Gour
}
\affiliations{
    Carnegie Mellon University\\
    5000 Forbes Avenue\\
    Pittsburgh, Pennsylvania 15213-3890 USA\\
    \{riyaza,shivansg\}@andrew.cmu.edu
}


\begin{document}

\maketitle

\begin{abstract}

\end{abstract}

\section{Introduction}
\subsection{Motivation} % problem statement and motivation
Modern-day scientific research, especially in verifiable fields such as mathematics and computer science, is highly inefficient. Researchers frequently duplicate efforts, pursue dead ends, and lack a proper mechanism for determining the true value of a research direction. With the advent of autonomous AI scientists, we risk aggravating these inefficiencies if such systems are improperly integrated into the current research landscape.

Currently, the AI industry relies on centralized model routers to allocate queries to pools of heterogeneous agents; however, mathematical discovery is too complex a task for these routers to efficiently allocate compute across. This is because of the fact that as tasks within mathematics become more complex, they require recursive breakdown into a myriad of subtasks and prerequisite lemmas and theorems. An external, centralized planner cannot perfectly map the heterogeneous capabilities of a high number of agents to such an intricate and deeply nested dependency structure, even with untenable amounts of compute. Moreover, agents understand their own proficiencies and computational limits much better than an external router.

To effectively allocate compute across the irregularities within the agents' capabilities, we ask how one may incentivize independent, self-interested agents to optimally distribute themselves across these subtasks. In this project, we propose a method to shift from centralized routing to a decentralized market, where we allow automated mathematicians to ``bid" on tasks with some compute budget and use financial incentives to drive both collaboration and competition among the agents. In doing so, we hope to align the rationality of each agent towards the efficiency of mathematical discovery.
\subsection{Overview} % agora and agora forge description
We first conceptualize \textit{Agora}, a research market game where multiple agents (LLMs) pursue interdependent mathematical targets (in Lean) and incentivize each other via tradable claims and bounties. Within this system, theorems are treated as financial assets, and agents invest compute to search for proofs and sell intermediate lemmas to the market, allowing the price signal to guide resource allocation.

However, testing Agora's market mechanism and equilibria with LLMs is both slow and computationally intensive. Inference latency makes it difficult to efficiently iterate on different market mechanisms, economic designs, or train agents with reinforcement learning (RL).

To bridge this gap, we present \textit{Agora Forge}, the primary focus of this report. Agora Forge is a synthetic, multiagent environment which formalizes mathematical discovery as a partially observable stochastic game. Instead of generating whole proofs in Lean or natural language, Agora Forge abstracts this process into a resource allocation problem over a directed acyclic graph of logical dependencies. By abstracting away the real-world complexities of formal theorem proving and reducing them to three core strategic primitives (conjecturing, proving, and publishing), Agora Forge provides a controlled, sandbox-like environment that allows for rapid prototyping of market mechanisms and training agents to optimize the system's welfare before deploying them in the full, real-world setting of Agora.
\section{Related Work}
\subsection{Math Markets}
\subsection{Prediction Markets}
\subsection{Routing and Orchestration}
\subsection{Multiagent Reinforcement Learning}
\subsection{Mechanism Design}



\section{Background}
\subsection{Preliminary Definitions}

\subsubsection{Syntax}

We consider the \textit{formula} $\phi$ as the core syntactic primitive of (formal) mathematics. Namely, we represent mathematical theorem statements syntactically as a formula (i.e. object of type \texttt{Prop}) within some formal language $\Sigma$, such as Lean, Rocq, etc. We define $\mathcal{F}$ as the collection of all such formulas, and equip it with an involution $\lnot : \mathcal{F}\rightarrow \mathcal{F}$; namely, $\lnot(\lnot\varphi)=\varphi$ for all $\varphi\in\mathcal{F}$. Note that we will refer to formulas as ``nodes'' interchangably due to a natural graphical structure that will be discussed later.

\subsubsection{Semantics}

We now aim to represent the semantic meaning of a theorem statement within a particular context concisely. Towards this, we consider the \textit{sequent}, given as an element of $\mathcal{S} := \mathcal{P}_{fin}(\mathcal{F}) \times \mathcal{F} $ and stylized as $[\Gamma \proves \phi]$ (read: ``Gamma \textit{proves} phi''). Namely, such a sequent represents a \textit{proof obligation}, i.e. the notion of proving a theorem statement $\phi$ (the \textit{consequent}) within the \textit{antecedent} context $\Gamma$. Such a proof obligation may be \textit{resolved}, meaning proven.

\subsubsection{Libraries}

We assume the existence of a universal clock $t \in \mathbb N$ and set of agents $\mathcal N$, and use it to define a library $\mathcal{L}_\bullet = \{\mathcal{L}_t\}_{t \in \mathbb{N}}$ as a collection of discrete library states that evolves over time. At every discrete timestep $t \in \mathbb{N}$, each library state is given as a tuple 
\[\mathcal{L}_t := (\mathcal{C}_t,\mathcal{RS}_t, \delta_t, \Theta_t) \in \mathscr{L}\]
where we partition $\mathcal{F} = \mathcal{C}_t \cup \mathcal{G}_t, \mathcal{C}_t \cap \mathcal{G}_t = \varnothing$ into \textit{concrete} and \textit{ghost} nodes. Intuitively, concrete nodes represent theorem statements that have been stated or conjectured, and ghost nodes represent all hypothetical, possible theorem statements that are yet to be conjectured or stated. $\mathcal{C}_t$ naturally induces a semantic analogue: let $\mathcal{CS}_t := \mathcal{P}_{fin}(\mathcal{C}_t) \times \mathcal{C}_t$ be the set of all \textit{concrete sequents}, which contains a subset $\mathcal{RS}_t \subseteq \mathcal{CS}_t$ of \textit{resolved sequents}, which represent the concrete nodes that have been proven (within a particular antecedent context, of course). 


Mathematically, a proof of a particular theorem in some context often explicitly depends on other lemmas within this context. In an effort to model this phenomenon, define $\delta_t : \mathcal{RS}_t \to \mathcal{P}_{fin}(\mathcal{F})$ as the dependency mapping.
Note that in mathematics, there may exist many proofs of the same theorem with different dependencies for each, and one may find new proofs of the same theorem that are also valid. Thus, in the most general formulation, we only require that resolved sequents must remain resolved ($\mathcal{RS}_{t} \subseteq \mathcal{RS}_{t+1}$), though the underlying ``proof'' (moreso, the set of dependencies of this ``proof'') can change.

And lastly, we define $\Theta_t : \mathcal{RS}_t\to\mathbb N^+ \times \mathcal N$ to record the (discrete) solve-time associated to each resolved sequent, and the agent responsible for the resolution. 
% as in our formulation we do not allow the modification or deletion of a proof after it is initially proven (as $\mathcal{RS}_{t_0} \subseteq \mathcal{RS}_{t}$ for $t_0 < t$).


So, in other words, the library state $\mathcal{L}_t$ represents the collection of all stated theorem statements, the collection of proof witnesses within particular contexts for a subset of these statements, and the dependencies and ancillary metadata of these witnesses.

\paragraph{Library Invariants and Properties}

Naturally, we impose a monotonicity invariant for libraries: one may not ``unconjecture'' or ``unprove'' a theorem. Formally, given a library $\mathcal{L}_\bullet = \{\mathcal{L}_t\}_{t \in \mathbb{N}}=\{(\mathcal{C}_t,\mathcal{RS}_t,\delta_t,\Theta_t)\}_{t \in \mathbb{N}}$, we require that for all $t \in \mathbb{N}$:
\[\mathcal{C}_{t} \subseteq \mathcal{C}_{t+1} \quad \text{and} \quad \mathcal{RS}_{t} \subseteq \mathcal{RS}_{t+1}\]

Additionally, we require that if one conjectures a theorem, they must also have automatically conjectured its converse; in other words:
\[\phi \in \mathcal{C}_t \iff \neg\phi \in \mathcal{C}_t\]

Similarly, one may not prove both a theorem and its converse true: 
\[[\Gamma \proves \phi] \in \mathcal{RS}_t \implies [\Gamma \proves \neg \phi] \notin \mathcal{RS}_t\]

Semantically, we note reflexivity holds: $\phi \in \Gamma \implies [\Gamma \proves \phi] \in \mathcal{RS}_t$. And for the sake of consistency, we require that for any $[\Gamma \proves \phi] \in \mathcal{CS}_t$, $\psi \in \Gamma \implies \neg\psi \notin \Gamma$ (although we do not impose any full propositional semantics or richer logical theory behind these abstract formulas, preferring to leave them fully symbolic objects with various resolution and concreteness states, allowing inconsistent antecedents in our witnesses limits the scalability of our approach to deployment as a strategy module for multiagent theorem proving in a fully-defined formal language and type theory). 

Note that this last condition implicitly redefines
\[\mathcal{CS}_t := \mathcal{P}_{fin}(\mathcal{C}_t)/\sim \times \mathcal{C}_t \] where $\phi\sim\neg\phi$ is the equivalence relation generated by the negation operation. So, due to this consistency condition, we note that $\mathcal{RS}_t$ is upwards-closed, in that if $[\Gamma \proves \phi] \in \mathcal{RS}_t$, then for all $\Gamma' \in \mathcal{P}_{fin}(\mathcal{C}_t)$, $\Gamma \subseteq \Gamma' \implies [\Gamma'\proves\phi]\in \mathcal{RS}_t$.

Additionally, considering the dependency mapping, we note that intuitively $\delta_t$ must satisfy the property that $\delta_t([\Gamma \proves \phi]) \subseteq \Gamma$. Moreover, one may use $\delta_t$ to induce a directed graph $D(\delta_t) := (\mathcal{CS}_t,\mathcal{E}_t)$, where 
\[\mathcal{E}_t = \Big\{([\Gamma\proves \phi], [\Gamma' \proves \psi]) : \phi \in \delta_t([\Gamma' \proves \psi]), \Gamma\subseteq \Gamma'\Big\}\]
We require that this induced directed graph $D(\delta_t)$ of $\delta_t$ must be acyclic.

\paragraph{Axioms and Fully Resolved Sequents}

As an aside, we additionally note that although $[\Gamma \proves \phi] \in \mathcal{RS}_t$ may be marked resolved, that does not mean that $\phi$ has a ``error''-free proof within the current library state as its proof may depend on some still-unresolved lemmas. More concretely, we may define $\mathcal{FRS}_t \subseteq \mathcal{RS}_t$ as the collection of \textit{fully resolved sequents}, which is defined inductively such that \textit{axioms} $[\proves \phi] \in \mathcal{FRS}_t$ and $[\Gamma \proves \phi] \in \mathcal{FRS}_t$ if and only if for all $\phi' \in \delta_t([\Gamma \proves \phi])$, there exists $\Gamma'\subseteq \Gamma$ such that $[\Gamma' \proves \phi'] \in \mathcal{FRS}_t$.


\paragraph{Private and Shared Libraries}
In the multiagent setting, each agent $\alpha\in\mathcal{N}$ maintains a \textit{private library} $\mathcal{L}_\bullet^\alpha = \{\mathcal{L}^\alpha_t\}_{t\in\mathbb{N}}$ with
\[
\mathcal{L}^\alpha_t := (\mathcal{C}^\alpha_t,\mathcal{RS}^\alpha_t,\delta^\alpha_t,\Theta_t^\alpha),
\]
and the environment maintains a \textit{shared (public) library} $\mathcal{L}_\bullet^0 = \{\mathcal{L}^0_t\}_{t\in\mathbb{N}}$ with
\[
\mathcal{L}^0_t := (\mathcal{C}^0_t,\mathcal{RS}^0_t,\delta^0_t, \Theta_t^0).
\]

We additionally enforce the following inclusion invariant for all $t \in \mathbb{N}, \alpha \in \mathcal{N}$:
\[
\mathcal{L}^0_t \subseteq \mathcal{L}^\alpha_t
\]
meaning $\mathcal{C}^0_t\subseteq \mathcal{C}^\alpha_t$, $\mathcal{RS}^0_t\subseteq \mathcal{RS}^\alpha_t$, and for all $s\in \mathcal{RS}_t^0, \delta^0_t(s) = \delta^\alpha_t(s)$ and $\Theta^0_t(s)=\Theta^\alpha_t(s)$.


\subsubsection{Market State}

We additionally model the market as a discrete-time process $\mathcal{M}_\bullet := \{\mathcal{M}_t\}_{t\in\mathbb{N}}$.
At each timestep $t$, the market state is
\[
\mathcal{M}_t := \big(\{\mathcal{P}^\alpha_t\}_{\alpha\in\mathcal{N}},\;\mathcal{ML}_t\big)
\in \mathfrak{P}^n \times \mathfrak{L} = \mathfrak{M}\]
where $\mathcal{P}^\alpha_t\in\mathfrak{P}$ is agent $\alpha$'s private portfolio state and $\mathcal{ML}_t \in \mathfrak{L}$
is the public market ledger.

We assume a canonical balance/valuation map $\beta:\mathfrak{P}\to\mathbb{R}$ and define the balance of agent $\alpha$ at time $t$ by
\[
B^\alpha_t := \beta(\mathcal{P}^\alpha_t)\in\mathbb{R}.
\]



\subsection{Action Primitives}

We now specify the core action primitives of VAMP. These primitives will later be assembled into the
POSG transition and observation structure; for now, we define them as standalone kernels and update maps.

\paragraph{Discrete kernels.}
All stochastic primitives in this section have finite or countable output spaces, and may be treated as
probability mass functions (pmfs). Concretely, for an input space $X$ and discrete output space $Y$,
a stochastic kernel is a map $K:Y\times X\to[0,1]$ such that $\sum_{y\in Y}K(y\mid x)=1$ for all $x\in X$.

\subsubsection{Proof Kernel}

Fix an agent $\alpha\in\mathcal{N}$. The proof kernel models allocating a time budget to attempt to resolve a target sequent.

% \paragraph{Type.}
Let the input space be
\[
X_{\mathrm{proof}} := \mathscr{L}\times\mathcal{S}\times \mathbb{N}^+,
\]
where $(s,\tau)$ consists of a target sequent $s=[\Gamma\proves\phi]$ and a discrete time budget $\tau$.
Let the output space be
\[
Y_{\mathrm{proof}} := \{(0,\bot)\}\;\cup\;\{(1,d): d\in\mathcal{P}_{fin}(\mathcal{F})\},
\]
where $(0,\bot)$ denotes failure, and $(1,d)$ denotes success together with a dependency set $d$.

The proof kernel for agent $\alpha$ is a pmf

\[\mathcal{K}^{\alpha}_{\mathrm{proof}} : Y_{\mathrm{proof}}\times X_{\mathrm{proof}} \to [0,1]\]


% \paragraph{Intended admissibility.}
Although $\mathcal{K}^{\alpha}_{\mathrm{proof}}$ is defined on all of $\mathcal{S}\times\mathbb{N}^+$,
in the intended game dynamics the agent is only allowed to invoke a proof attempt on \emph{concrete} targets
available in its local library (later captured via a state-dependent admissibility predicate).
On admissible inputs $(\mathcal{L},[\Gamma\proves\phi],\tau)$, a successful output $(1,d)$ is intended to satisfy $d\subseteq\Gamma$.

% \paragraph{Intended state effect (informal).}
Upon success, the environment will (i) mark the target sequent resolved in $\alpha$'s library and
(ii) record/update its dependency map at that sequent with the returned $d$, subject to global consistency
(e.g., preserving acyclicity of the induced dependency graph). Upon failure, no new resolved sequent is added.
These effects are enforced later by the global state-update operator.

\subsubsection{Conjecture Kernel}

Fix an agent $\alpha\in\mathcal{N}$. The conjecture kernel models allocating a time budget to propose a new formula,
anchored at an existing sequent.

Let the input space be
\[
X_{\mathrm{conj}} :=\mathscr{L}\times \mathcal{S}\times \mathbb{N}^+,
\]
where $(s,\tau)$ consists of an anchor sequent $s=[\Gamma\proves\phi]$ and time budget $\tau$.
Let the output space be
\[
Y_{\mathrm{conj}} := \{(0,\bot)\}\;\cup\;\{(1,\psi): \psi\in\mathcal{F}\},
\]
where $(0,\bot)$ denotes failure and $(1,\psi)$ denotes successful conjecture of formula $\psi$.

The conjecture kernel for agent $\alpha$ is a pmf
\[
\mathcal{K}^{\alpha}_{\mathrm{conj}} : Y_{\mathrm{conj}}\times X_{\mathrm{conj}} \to [0,1].
\]

As with the proof kernel, $\mathcal{K}^{\alpha}_{\mathrm{conj}}$ is only intended to be invoked on anchors available to the agent (e.g., concrete sequents in $\alpha$'s local library).

Moreover, given admissible input $(\mathcal L, [\Gamma \proves \phi], \tau)$ and a successful output $(1,\psi) \sim \mathcal{K}_{\mathrm{conj}}^\alpha (\cdot \mid \mathcal L, [\Gamma \proves \phi], \tau)$, we enforce that $\psi$ is a ghost node. Namely, if $\mathcal{L} = (\mathcal{C},\mathcal{RS},\delta,\Theta)$, then $\psi \in \mathcal{F} \setminus \mathcal{C}$.


Upon such success with output $(1,\psi)$, we intend the environment to add $\psi$ to $\alpha$'s set of concrete formulas (and also add $\neg\psi$ to preserve the negation-pair invariant). This converts $\psi$ from ghost to concrete relative to $\alpha$.

\subsubsection{Dependency Closure for Publication}

Publication is a deterministic primitive that copies a closure-completed subset of an agent's private library
into the shared library, and (by the shared-library invariant) into every agent's local library.

Fix agent $\alpha$ and time $t$, and suppose $\delta_t^\alpha$ induces the directed acyclic graph
$D(\delta_t^\alpha)=(\mathcal{CS}_t^\alpha,\mathcal{E}_t^\alpha)$ as defined earlier.
For any set of target resolved sequents $R\subseteq \mathcal{RS}_t^\alpha$, define its \emph{dependency closure}:
\begin{align*}
\mathrm{cl}^{\alpha}_t(R) := 
\Big\{\, &s\in \mathcal{RS}_t^\alpha : \exists r\in R \text{ such that there}\\& \text{is a directed path } s\to r
\text{ in } D(\delta_t^\alpha)\,\Big\}
\end{align*}

Equivalently, $\mathrm{cl}^\alpha_t(R)$ is the set of all resolved sequents required (transitively) as dependencies
of the sequents in $R$.


We now proceed to extend this closure operator to the payload for such publication actions. Namely, a publication payload is a pair $U_t^\alpha=(C,R)$ with $C\subseteq \mathcal{C}_t^\alpha,\;\; R\subseteq \mathcal{RS}_t^\alpha$.

Define its closure-completion $\overline{U_t^\alpha}:=\big(\overline{C},\overline{R}\big)$ where $\overline{R}:=\mathrm{cl}^\alpha_t(R)$ and
\[\overline{C}:= C \;\cup\; \bigcup_{[\Gamma\proves\phi]\in \overline{R}}(\Gamma\cup\{\phi\})\]

Thus, publishing resolved sequents forces publication of all formulas appearing in the antecedents and consequents of the dependency-closed witness set.


Therefore, given $\overline{U_t^\alpha}$, intuitively, the environment updates the shared library by unioning in $\overline{C}$ and $\overline{R}$ (and any associated dependency information it chooses to expose),
then propagates these additions to every agent's local library so that $\mathcal{L}^0_{t+1}\subseteq \mathcal{L}^\beta_{t+1}$
continues to hold for all $\beta$.



\subsubsection{Query Model}

Each agent $\alpha$ maintains a stateful \emph{query model} which estimates how its prover stack will perform on a
requested proof search. Intuitively, the query model is a continually-updated predictor for
(i) the probability of successfully resolving a target sequent within a given time budget and
(ii) the expected time-to-solution (or a proxy thereof).


Fix an abstract space $\mathcal{Q}$ of query models (e.g., parameters of a Bayesian estimator, weights of a learned predictor, or a table of sufficient statistics). At each timestep $t$, agent $\alpha$ has a current query model $Q_t^\alpha \in \mathcal{Q}$.
A query operation takes as input a target (concrete) sequent:
\[
q^\alpha : \mathcal{Q}\times \mathcal{S} \to [0,1]\times \mathbb{N}
\]
and write
\[
(\hat p_t^\alpha(s),\hat \tau_t^\alpha(s))
:=
q^\alpha(Q_t^\alpha,s),
\]
where $\hat p_t^\alpha([\Gamma \proves \phi])$ is the model's estimate of $\Pr(\phi\text{ is true})$
and $\hat\tau_t^\alpha(s)$ is an estimate of the expected budget required to successfully prove $[\Gamma\proves \phi]$. 

More explicitly, for a fixed agent $\alpha$ and library state $\mathcal{L}$, and target sequent $s\in\mathcal{S}$, define the probability of success within budget $\tau$ by
\[
p^\alpha_\tau(s) \;:=\; \sum_{d\in\mathcal{P}_{fin}(\mathcal{F})}\mathcal{K}^{\alpha}_{\mathrm{proof}}\big((1,d)\mid \mathcal L, s,\tau\big).
\]
This induces a time-to-solution distribution by declaring a random variable $T^\alpha_s$ with CDF
\[
\Pr\!\big(T^\alpha_s \le \tau\big) \;:=\; p^\alpha_\tau(s), \qquad \tau\in\mathbb{N}^+.
\]
We interpret the eventual solvability (or ``truth/resolvability'' under $\alpha$'s prover stack) as
\[
p^\alpha_\infty(s) \;:=\; \sup_{\tau\in\mathbb{N}^+} p^\alpha_\tau(s),
\]
and, when $p^\alpha_\infty(s)>0$, the conditional expected solve time as
\[
\tau^\alpha_\infty(s) \;:=\; \mathbb{E}\!\left[T^\alpha_s \,\middle|\, T^\alpha_s < \infty\right]=\sum_{k\ge1}\Pr(T^\alpha_s\ge k \mid T^\alpha_s<\infty)
\]
Accordingly, the query readout $q^\alpha(Q^\alpha_t,s)=(\hat p^\alpha_t(s),\hat\tau^\alpha_t(s))$ is intended to approximate
\[
\hat p^\alpha_t(s)\approx p^\alpha_\infty(s),
\qquad
\hat\tau^\alpha_t(s)\approx \tau^\alpha_\infty(s),
\]
using the agent's accumulated evidence.



The query model is updated from observed outcomes of proof attempts (and optionally from public disclosures).

Define the space of query-training data as
\[
\mathcal{D}(\mathcal{Q}) := \mathcal{S} \times (\mathbb{N}^+ \times \mathcal{N}),
\]
where a datum $(s,(\tau, \alpha'))$ denotes a successful proof attempt by agent $\alpha'$ on $s$ on time budget $\tau$.

Define an update operator
\[
\Delta^\alpha : \mathcal{Q}\times \mathcal{D}(\mathcal{Q}) \to \mathcal{Q}
\]

Given an observed training datum $d_t^\alpha\in\mathcal{D}(\mathcal{Q})$ available to agent $\alpha$ at time $t$, the query model updates as
\[
Q_{t+1}^\alpha :=\Delta^\alpha(Q_t^\alpha, d_t^\alpha),
\]
We enforce the commutativity and associativity of $\Delta^\alpha$, so that if multiple data points $\{(d_t^\alpha)_i\}_{i=0}^{k-1}$ are observed in a single timestep, there is a single and well-defined batch update which is the composition (in any order) of the individual datum updates: i.e., if $Q_t^\alpha=(Q_t^\alpha)_0$,  
\[
(Q_{t}^\alpha)_{i+1} = \Delta^\alpha((Q_t^\alpha)_i, (d_t^\alpha)_i)
\]
\[Q_{t+1}^\alpha := (Q_{t}^\alpha)_{k-1}\]
which is consistent across all possible permutations on $\{0,\cdots,k-1\}$.



The framework does not prescribe the internal form of $Q_t^\alpha$, and allows downstream experiments to arbitrarily instantiate $(\mathcal{Q}, q^\alpha, \Delta^\alpha)$ as desired, whether that means querying based on Bayesian update rules, learned predictors, or any other mechanism.

\paragraph{Market Mechanism}
\label{sec:vamp-market}

We briefly take an aside to define the \emph{market mechanism} as a modular plugin that that governs (i) which market actions are available, (ii) which are admissible at a given world state, and (iii) how the market state evolves. One may parameterize and study many different verifiable allocation games by simply switching out this underlying economic reward mechanism.
\renewcommand{\epsilon}{\varepsilon}
We define a market mechanism $\mathscr{M}$ as a tuple:
\begin{align*}
\mathscr{M} :=\Big(&\mathcal{N}, \mathfrak{M},\beta, \mathcal{W},\{\mathcal{A}^\alpha_{\mathrm{mkt}}\}_{\alpha\in\mathcal{N}},
\\&\{\mathrm{Adm}^\alpha_{\mathrm{mkt}}\}_{\alpha\in\mathcal{N}},\mathcal{T}_{\mathrm{mkt}}\Big)
\end{align*}
where:
\begin{itemize}
    \item $\mathcal{N}$ is the set of agents, indexed by $\alpha\in \mathcal{N}$.
    \item $\mathfrak{M}=(\mathfrak{P}^n,\mathfrak{L})$ is the market state space, where $\mathfrak{P}$ is the space of possible portfolios and $\mathfrak{L}$ is the space of possible public market ledgers.
    \item $\beta : \mathfrak{P} \to \mathbb{R}$ is the canonical balance/valuation map.
    \item $\mathcal{W}$ is the world (state) space. Note that $\mathfrak{M}$ is a component in $\mathcal{W}$.
  
    \item $\mathcal{A}^\alpha_{\mathrm{mkt}}$ is the (ambient) \textit{market} action space of agent $\alpha$. Note that the market no-op is always a valid action: $\bot_{\mathrm{mkt}} \in \mathcal{A}_{\mathrm{mkt}}^\alpha$. The joint market action space is:
    \[
\mathcal{A}_{\mathrm{mkt}} := \prod_{\alpha\in\mathcal{N}} \mathcal{A}^\alpha_{\mathrm{mkt}}
\]
    \item $\mathrm{Adm}^\alpha_{\mathrm{mkt}}:\mathcal{W}\times\mathcal{A}^\alpha_{\mathrm{mkt}}\to\{0,1\}$ is the market admissibility predicate for agent $\alpha$. It induces the state-dependent admissible market action set:
    \[
    \mathcal{A}^\alpha_{\mathrm{mkt}}(s) := \{a^\alpha\in\mathcal{A}_{\mathrm{mkt}}^\alpha : \mathrm{Adm}_{\mathrm{mkt}}^\alpha(s,a^\alpha)=1\}.
    \]
    \item $\mathcal{T}_{\mathrm{mkt}} : \Big(\mathcal{W} \times \mathcal{W} \times \mathcal{A}_{\mathrm{mkt}}\Big) \times \Sigma_{\mathcal{W}} \to [0,1]$ is a stochastic transition kernel with input space $\mathcal{W} \times \mathcal{W} \times \mathcal{A}_{\mathrm{mkt}}$ and output space $\mathcal{W}$. Namely, given a change in world state from $w \to w' \in \mathcal{W}$, and a joint market action $a_{\mathrm{mkt}}\in\mathcal{A}_{\mathrm{mkt}}$, this kernel induces a distribution over the next world state:
\[
\bar{w}\sim \mathcal{T}_{\mathrm{mkt}}(\cdot\mid w,w',a_{\mathrm{mkt}}).
\]
Deterministic mechanisms are recovered as Dirac kernels.

    
    
    
\end{itemize}



\subsection{Partially Observable Stochastic Games}

A \emph{partially observable stochastic game} (POSG) is a mathematical model for sequential decision-making with multiple interacting agents under partial information. In this work we focus on the standard \emph{infinite-horizon discounted} formulation.
\subsubsection{Definition}

An infinite-horizon discounted POSG is specified by the tuple
\[
\Big(\mathcal{N}, \mathcal{W},\{\mathcal{O}^\alpha\}_{\alpha\in \mathcal{N}},\{\mathcal{Z}^\alpha\}_{\alpha\in \mathcal{N}},
\{\mathcal{A}^\alpha\}_{\alpha\in \mathcal{N}},\]\[\{\mathrm{Adm}^\alpha\}_{\alpha\in\mathcal{N}},\mathcal{T},
\{r^\alpha\}_{\alpha\in \mathcal{N}},\gamma,b_0\Big)    
\]

where:
\begin{itemize}
    \item $\mathcal{N}$ is the set of agents, indexed by $\alpha\in \mathcal{N}$.
    \item $\mathcal{W}$ is the world (state) space.
    \item $\mathcal{A}^\alpha$ is the (ambient) action space of agent $\alpha$, and the joint action space is
    $\mathcal{A}:=\prod_{\alpha\in \mathcal{N}}\mathcal{A}^\alpha$.
    \item $\mathcal{O}^\alpha$ is the observation space of agent $\alpha$.
    \item $\mathrm{Adm}^\alpha:\mathcal{W}\times \mathcal{A}^\alpha\to\{0,1\}$ is the \emph{admissibility predicate} for agent $\alpha$.
    It induces the state-dependent legal action set
    \[
    \mathcal{A}^\alpha(s) := \{a^\alpha\in\mathcal{A}^\alpha : \mathrm{Adm}^\alpha(s,a^\alpha)=1\}.
    \]
    We assume $\mathcal{A}^\alpha(s)\neq\varnothing$ for all $s$ (e.g.\ containing a no-op action).
    \item $\mathcal{T}$ is the transition kernel: for $s,s'\in \mathcal{W}$ and joint action $a\in\mathcal{A}$,
    \[
    \mathcal{T}(s' \mid s,a)\in[0,1]
    \]
    is the probability of transitioning from $s$ to $s'$ under $a$.
    \item $\mathcal{Z}^\alpha$ is the observation kernel for agent $\alpha$: for $o^\alpha\in\mathcal{O}^\alpha$,
    \[
    \mathcal{Z}^\alpha(o^\alpha \mid s,a,s')\in[0,1]
    \]
    is the probability that agent $\alpha$ observes $o^\alpha$ given transition $s\xrightarrow{a}s'$.
    \item $r^\alpha:\mathcal{W}\times \mathcal{A}\to \mathbb{R}$ is the per-step reward function for agent $\alpha$.
    \item $\gamma\in[0,1)$ is the discount factor.
    \item $b_0$ is the initial distribution over states $s_0\sim b_0$.
\end{itemize}

At each timestep $t$, each agent selects an action $a_t^\alpha\in \mathcal{A}^\alpha(s_t)$, yielding a joint action
$a_t=(a_t^\alpha)_{\alpha\in\mathcal{N}}\in\mathcal{A}$.
(Equivalently, one may allow agents to propose any $a^\alpha\in\mathcal{A}^\alpha$ and define a convention for inadmissible actions,
such as rejection/no-op or penalties; we use explicit admissibility predicates for clarity.)

\subsubsection{Policies, Utilities, and Equilibria}

A (behavioral) policy for agent $\alpha$ is a mapping from the player's action--observation history to a distribution over actions.
Let
\[
h^\alpha_t := (o^\alpha_0,a^\alpha_0,o^\alpha_1,a^\alpha_1,\dots,o^\alpha_t)\in(\mathcal{O}^\alpha\times \mathcal{A}^\alpha)^t\times \mathcal{O}^\alpha
\]
denote agent $\alpha$'s history up to time $t$. A policy $\pi^\alpha(\cdot \mid h^\alpha_t)$ is a distribution over $\mathcal{A}^\alpha$ conditioned on the past history $h^\alpha_t$. We write $\Pi^\alpha$ for the space of policies available to agent $\alpha$.

A \emph{pure strategy} is a map $\sigma^\alpha$ that selects a single action deterministically as a function of $h^\alpha_t$. A joint policy is $\pi := (\pi^\alpha)_{\alpha\in \mathcal{N}}$, and we write $\pi^{-\alpha}$ for the policies of all players except $\alpha$.

Given a joint policy $\pi$, player $\alpha$'s expected discounted return is
\[
R^\alpha(\pi) := \mathbb{E}_{\mathcal{T},\mathcal{Z},\pi}\!\left[\sum_{t=0}^{\infty}\gamma^t\, r^\alpha(s_t,a_t)\;\middle|\; s_0\sim b_0\right]
\]
where the expectation is taken over the the transition kernels $T$, observation kernels $Z$, and the (possibly stochastic) action selection induced by $\pi$. Under bounded rewards (e.g.\ $|r^\alpha(s,a)|\le R_{\max}$) and $\gamma\in[0,1)$, this sum converges almost surely
and $|R^\alpha(\pi)|\le \frac{R_{\max}}{1-\gamma}$.

Because $R^\alpha(\pi)$ depends on the joint policy and reward functions need not align across players, a POSG is not generally a single-agent optimization problem. Rather, one studies various behavioral policies and game-theoretic equilibria. One such equilibrium is the \emph{Nash equilibrium}: a joint policy $\pi$ is a Nash equilibrium if every agent $\alpha$ plays a best response to $\pi^{-\alpha}$, i.e., for all policies $\pi^{\alpha'}\in \Pi^\alpha$,
\[
R^\alpha(\pi^\alpha,\pi^{-\alpha}) \;\ge\; R^{\alpha}(\pi^{\alpha'},\pi^{-\alpha})
\]

% This is a rather strong condition, and one may also want to study weaker notions of equilibrium behaviors, such as correlated equilibria or coarse-correlated equilibria. \textbf{TODO DEFINE!}

A particularly important taxonomic feature of POSGs is the relationship between agents' reward functions. In \emph{cooperative} games, all agents share the same reward function: $r^\alpha(s,a)=r^{\alpha'}(s,a)$ for all $i,j$. In a two-player \emph{zero-sum} game,
$r^0(s,a)=-r^1(s,a)$ for all $(s,a)$. When no such restriction is imposed, the POSG is termed \emph{general-sum}.



\section{Verifiable Allocation Market POSG (VAMP)}
\label{sec:vamp}

We now define the \emph{Verifiable Allocation Market POSG} (VAMP) as a particular POSG designed to model multiagent verifiable subtask allocation mediated by an arbitrary market mechanism -- interpreted as an abstracted decentralized formal theorem proving environment.

Fix a set of agents $\mathcal{N}$. A VAMP is a class of infinite-horizon discounted POSG's, parameterized by a particular market mechanism $\mathscr{M}$:
\[
\Big(\mathcal{N},\mathcal{W},\{\mathcal{O}^\alpha\}_{\alpha\in\mathcal{N}},\{\mathcal{Z}^\alpha\}_{\alpha\in\mathcal{N}},\{\mathcal{A}^\alpha\}_{\alpha\in\mathcal{N}},
\]\[\{\mathrm{Adm}^\alpha\}_{\alpha\in\mathcal{N}},\mathcal{T},
\{r^\alpha\}_{\alpha\in\mathcal{N}},\gamma,b_0, \mathscr{M}\Big)
\]
defined as follows.

\subsection{World State Space $\mathcal{W}$}

A world state at time $t$ is a tuple
\begin{align*}
w_t :=
\Big(
&\{\mathcal{L}^\alpha_t\}_{\alpha\in\mathcal{N}},
\ \mathcal{L}^0_t,
\ \mathcal{M}_t,
\ \{J^\alpha_t\}_{\alpha\in\mathcal{N}},
\ \{Q^\alpha_t\}_{\alpha\in\mathcal{N}},
\\&\ \{\Omega^\alpha_t\}_{\alpha\in\mathcal{N}},
\ \{U_{t,\mathrm{proof}}^\alpha\}_{\alpha\in\mathcal N},
\ \{U_{t,\mathrm{conj}}^\alpha\}_{\alpha\in\mathcal N}
\Big)\in\mathcal{W},
\end{align*}
where:
\begin{itemize}
\item $\mathcal{L}^\alpha_t$ is agent $\alpha$'s private library and
$\mathcal{L}^0_t$ is the shared public library.
\item $\mathcal{M}_t=(\{\mathcal{P}^\alpha_t\}_{\alpha\in\mathcal{N}},\;\mathcal{ML}_t) \in \mathfrak{P}^n \times \mathfrak{L} =\mathfrak{M}$ is the market state, as defined in $\mathscr{M}$.
\item $J^\alpha_t$ is the durative job descriptor
\[
J^\alpha_t \in \{\bot\}\ \cup\ \big(\{\mathrm{Proof}, \mathrm{Conj}\}\times\mathcal{S}\times\mathbb{N}^+\times\mathbb{N}^+\big)
\]
with $\bot$ meaning idle and $(\cdot,s,\tau_{\mathrm{rem}},\tau_{\mathrm{eff}})$ meaning:
(i) $\tau_{\mathrm{rem}}$ steps remain in the \emph{current run} and
(ii) $\tau_{\mathrm{eff}}$ is the \emph{effective cumulative budget} used to evaluate the run upon completion.
Concretely, when a run is initiated at time $t$ with intended additional duration $\tau_{\mathrm{dur}}$ on target $s$,
we set $\tau_{\mathrm{eff}} := U^\alpha_{t,\mathrm{-}}(s) + \tau_{\mathrm{dur}}$ and $\tau_{\mathrm{rem}} := \tau_{\mathrm{dur}}$.
The agent is forced to take $\tau_{\mathrm{dur}}$ consecutive no-op actions, and at completion the outcome is sampled from the
appropriate kernel using budget $\tau_{\mathrm{eff}}$.
\item $Q^\alpha_t\in\mathcal{Q}$ is agent $\alpha$'s query model state, updated by $\Delta^\alpha$ on successful proofs.
\item  $\Omega^\alpha_t \in \{\bot\}\ \sqcup\ (\mathcal{S}\times[0,1]\times\mathbb{N})$ is agent $\alpha$'s last query response, either empty $\bot$ or a triple $(s,\hat p,\hat \tau)$ signifying an estimate that the sequent $s$ is true with probability $\hat p$, and given it is true, an estimate of the time it would take to prove it $\hat \tau$.
\item $U^\alpha_{t,\mathrm{-}} : \mathcal S \to \mathbb N$ records the total time already spent by agent $\alpha \in \mathcal N$ attempting a proof or conjecture for sequent $s \in \mathcal S$, up to time $t$. This is used to allow agents to ``pick up where they left off''. Note $U^\alpha_{0,\mathrm{-}}\equiv 0$.
\end{itemize}
We take $\mathcal{W}$ to be the set of tuples satisfying all such tuples satisfying the consitituent objects' various invariants and conditions.


\subsection{Observations $\mathcal{O}^\alpha$ and Observation Kernels $\mathcal{Z}^\alpha$}

% Public events are defined as the change in the public projection of the world state:
% \[
% \mathrm{Pub}(w_t):=(\mathcal{L}^0_t,\Theta^0_t,\mathcal{ML}_t),
% \]\[E^{\mathrm{pub}}_t := \mathrm{Pub}(w_{t+1}) \setminus \mathrm{Pub}(w_{t}),
% \]
% Thus, any update to a public field is automatically publicly visible.

Agent $\alpha$ observes
\[
o^\alpha_t :=
\big(\mathcal{L}^\alpha_t,\mathcal{P}^\alpha_t,B^\alpha_t, \Omega_t^\alpha;\ \mathcal{L}^0_t,\mathcal{ML}_t\big)%;\ E^{\mathrm{pub}}_t\big)
\in\mathcal{O}^\alpha.
\]
The observation kernel is deterministic:
\[
\mathcal{Z}^\alpha(o^\alpha\mid w,a,w')=\mathbf{1}[\,o^\alpha=\mathrm{Obs}^\alpha(w)\,],
\]
where $\mathrm{Obs}^\alpha(w)$ computes the observation slice as above.


\subsection{Action Space $\mathcal{A}^\alpha$ and Admissibility $\mathrm{Adm}^\alpha$}

For each agent $\alpha$, define the tagged \textit{non-market} action space
\begin{align*}
\mathcal{A}_{\mathrm{nm}}^\alpha
:=
&
\{\mathrm{Prove}\}\times(\mathcal{S}\times\mathbb{N}^+)
\\&\sqcup
\{\mathrm{Conj}\}\times(\mathcal{S}\times\mathbb{N}^+)
\\&\sqcup
\{\mathrm{Pub}\}\times(\mathcal{P}_{fin}(\mathcal{F})\times \mathcal{P}_{fin}(\mathcal{S}))
\\&\sqcup
\{\mathrm{Qry}\}\times \mathcal{S}
\\&\sqcup\{\bot_{\mathrm{nm}}\}
\end{align*}
Where $\bot_{\mathrm{nm}}$ indicates a non-market no-op. Thus, the full action space for agent $\alpha$ may be written as:
\[
\mathcal{A}^\alpha = \mathcal{A}^\alpha_{\mathrm{nm}} \sqcup \mathcal{A}^\alpha_{\mathrm{mkt}}
\]
where $\mathcal{A}^\alpha_{\mathrm{mkt}}$ is specified in $\mathscr{M}$. Note then that as an agent $\alpha$ may only take one action at a given timestep, both $\bot_{\mathrm{nm}}$ and $\bot_{\mathrm{mkt}}$ would indicate a no-op.

As such, the joint action space is 
\[\mathcal{A}:=\prod_{\alpha\in\mathcal{N}}\mathcal{A}^\alpha\]


Note that for a  given joint action $a \in \mathcal{A}$, we may uniquely decompose it into a non-market and market joint action. In other words:
\[a = a_{\mathrm{nm}} \oplus a_{\mathrm{mkt}}\]
Where for each agent-component of the joint action $a^\alpha \in \mathcal{A}^\alpha$, if $a^\alpha \in \mathcal{A}^\alpha_{\mathrm{nm}}$, then let $a^\alpha_{\mathrm{nm}} = a^\alpha$ and $a^\alpha_{\mathrm{mkt}} = \bot_{\mathrm{mkt}}$. Alternatively, if $a^\alpha \in \mathcal{A}^\alpha_{\mathrm{mkt}}$, then let $a^\alpha_{\mathrm{mkt}} = a^\alpha$ and $a^\alpha_{\mathrm{nm}} = \bot_{\mathrm{nm}}$. Thus, we may define:
\[a_{\mathrm{nm}} = \prod_{\alpha \in \mathcal{N}} a^\alpha_{\mathrm{nm}},\qquad a_{\mathrm{mkt}} = \prod_{\alpha \in \mathcal{N}} a^\alpha_{\mathrm{mkt}}\]




Now, define $\mathrm{Adm}^\alpha:\mathcal{W}\times\mathcal{A}^\alpha\to\{0,1\}$ as an admissibility predicate, which enforces semantic and operational constraints. 

Fix $w_t \in \mathcal{W}$ and consider $\alpha \in \mathcal{N}$. Then:
\begin{itemize}
    \item If $J^\alpha_t=\bot$, then:
\begin{itemize}
\item $\mathrm{Adm}^\alpha(w_t,\bot)=1$.
\item For $a^\alpha_t = (\mathrm{Prove},(s,\tau))$ and $a^\alpha_t = (\mathrm{Conj},(s,\tau))$:
\[
\mathrm{Adm}^\alpha(w_t,a^\alpha_t)=1 \iff s \in \mathcal{CS}_t^\alpha
\]
\item For $a^\alpha_t = (\mathrm{Pub},(C,R))$: \[\mathrm{Adm}^\alpha(w_t,a^\alpha_t)=1 \iff C\subseteq\mathcal{C}_t^\alpha,\ R\subseteq\mathcal{RS}_t^\alpha\]
\item For $a^\alpha_t = (\mathrm{Qry},s)$: 
\[\mathrm{Adm}^\alpha(w_t,a_t^\alpha)=1 \iff s \in \mathcal{CS}_t^\alpha\]
\item For market actions $a^\alpha_t\in\mathcal{A}^\alpha_{\mathrm{mkt}}$:
\[
\mathrm{Adm}^\alpha(w_t,a^\alpha_t)= \mathrm{Adm}^\alpha_{\mathrm{mkt}}(w_t,a^\alpha_t) 
\]
where $\mathrm{Adm}_{\mathrm{mkt}}^\alpha$ is specified by $\mathscr{M}$.
\end{itemize}

\item If $J^\alpha\neq\bot$, then the agent is forced to take no-op: $\mathrm{Adm}^\alpha(w_t,\bot)=1$ and for all $a^\alpha_t\in \mathcal{A}^\alpha\setminus\{\bot\}$, $\mathrm{Adm}^\alpha(w_t,a^\alpha_t)=0$.
\end{itemize}

\subsection{Transition Kernel $\mathcal{T}$}

We define $\mathcal{T}(\cdot\mid w_t,a_t)$ by a two-stage generative construction. Namely, given an admissible input pair $(w_t,a_t) \in \mathcal{W} \times \mathcal{A}$, uniquely decompose the joint action as
\[
a_t = a_{t,\mathrm{nm}} \oplus a_{t,\mathrm{mkt}}
\]

We then define a non-market transisition kernel $\mathcal{T}_{\mathrm{nm}} : (\mathcal{W} \times \mathcal{A}_\mathrm{nm}) \times \Sigma_{\mathcal{W}} \to [0,1]$ and consider the market transition kernel $\mathcal{T}_{\mathrm{mkt}} : (\mathcal{W} \times \mathcal{W} \times \mathcal{A}_\mathrm{mkt}) \times \Sigma_{\mathcal{W}} \to [0,1]$ as specified by $\mathscr{M}$. Intuitively, the full transition kernel $\mathcal{T}$ is the composition of the two stages:
first draw $w'_t\sim \mathcal{T}_{\mathrm{nm}}(\cdot\mid w_t,a_{t,\mathrm{nm}})$, then draw
$w_{t+1}\sim \mathcal{T}_{\mathrm{mkt}}(\cdot\mid w_t,w'_t,a_{t,\mathrm{mkt}})$.

We now proceed to define each stage of this construction.

\subsubsection{Background: Deterministic Pre-Update}
 
 We introduce a deterministic map that serves as the first step of the non-market transition kernel. It applies everything that happens immediately from the joint non-market action at this step, before any stochastic proof/conjecture outcomes are realized. So, given a world state $w \in \mathcal{W}$, we define:
\[\tilde w := \mathrm{Pre}(w,\mathbf a_{\mathrm{nm}})\]
as the world state obtained by applying the following deterministic edits in order:
\begin{enumerate}
    \item Decrement job timers. Namely, for each $\alpha \in \mathcal{N}$:
    \begin{itemize}
        \item If $J^\alpha=\bot$, leave it.
	\item If $J^\alpha=(\mathrm{Proof},s,\tau_{\mathrm{rem}},\tau_{\mathrm{tot}})$ or $(\mathrm{Conj},s,\tau_{\mathrm{rem}},\tau_{\mathrm{tot}})$ with $\tau_{\mathrm{rem}}>0$, set
$\tau_{\mathrm{rem}}\leftarrow \tau_{\mathrm{rem}}-1$.
    \end{itemize}
\item Start new proof/conjecture jobs if idle. Namely, for each $\alpha \in \mathcal N$ with $J^\alpha=\bot$:
\begin{itemize}
    \item 	If $a^\alpha_{\mathrm{nm}}=(\mathrm{Prove},(s,\tau))$, set
$J^\alpha \leftarrow (\mathrm{Proof},s,\tau,\tau+U^\alpha_{\mathrm{proof}}(s))$.
	\item	If $a^\alpha_{\mathrm{nm}}=(\mathrm{Conj},(s,\tau))$, set
$J^\alpha \leftarrow (\mathrm{Conj},s,\tau,\tau+U^\alpha_{\mathrm{conj}}(s))$.
	\item	Otherwise do nothing.
\end{itemize}
\item Apply publication actions and propagate to shared library. Namely, for each $\alpha\in\mathcal{N}$, if $a^\alpha_{\mathrm{nm}}=(\mathrm{Pub},(C^\alpha,R^\alpha))$, then consider the payload $U^\alpha:=(C^\alpha,R^\alpha)$. Compute the closure $\overline{U^\alpha}$ and update the shared library and query model. Namely, let $\mathcal{C}^0_\mathrm{old} := \mathcal{C}^0$ and $\mathcal{RS}^0_\mathrm{old} := \mathcal{RS}^0$, and update:
\[\mathcal C^0 \leftarrow \mathcal C^0 \cup \bigcup_{\alpha}\overline C^\alpha,
\qquad
\mathcal{RS}^0 \leftarrow \mathcal{RS}^0 \cup \bigcup_{\alpha}\overline R^\alpha\]

Then for all $s \in \mathcal{RS}^0 \setminus \mathcal{RS}^0_\mathrm{old}$:
\[\delta^0(s) \leftarrow \delta^\alpha(s),\qquad  \Theta^0(s) \leftarrow \Theta^\alpha(s)\]

Then propagate public additions to each private library to maintain the invariant $\mathcal L^0\subseteq \mathcal L^\beta$:
\[\mathcal C^\beta \leftarrow \mathcal C^\beta \cup (\mathcal C^0\setminus \mathcal C^0_{\text{old}})\]\[
\mathcal{RS}^\beta \leftarrow \mathcal{RS}^\beta \cup (\mathcal{RS}^0\setminus \mathcal{RS}^0_{\text{old}})\]
Also update for all $s\in \mathcal{RS}^0 \setminus \mathcal{RS}^0_\mathrm{old}$:
 \[\delta^\beta(s)\leftarrow\delta^0(s),\qquad \Theta^\beta(s)\leftarrow\Theta^0(s)\]
 \[Q^\beta \leftarrow \Delta^\beta(Q^\beta,(s,\Theta^0(s)))\]
 


\item Update query buffer. Namely, for each $\alpha\in\mathcal{N}$, if $a^\alpha_{\mathrm{nm}}=(\mathrm{Qry},s)$, set
\[\Omega^\alpha \leftarrow \big(s,\hat p,\max\{0,\hat\tau-U^\alpha_{\mathrm{proof}}(s)\}\big)\]
where $(\hat p,\hat\tau):=q^\alpha(Q^\alpha,s)$. Otherwise set $\Omega^\alpha\leftarrow \bot$.

\end{enumerate}

\subsubsection{Background - Stochastic Job Completion}

Given $\tilde w$, define completion index sets:
\[
    \mathsf{Comp}_{\mathrm{proof}}(\tilde w) = \{\alpha : J^\alpha(\tilde w)=(\mathrm{Proof},s,0,\tau)\text{ for some }s,\tau\}
\]
\[\mathsf{Comp}_{\mathrm{conj}}(\tilde w) = \{\alpha : J^\alpha(\tilde w)=(\mathrm{Conj},s,0,\tau)\text{ for some }s,\tau\}\]

For each $\alpha\in \mathsf{Comp}_{\mathrm{proof}}(\tilde w)$, let $(s^\alpha_{\mathrm{proof}}(\tilde w),\tau^\alpha_{\mathrm{proof}}(\tilde w))$ be the $(s,\tau)$ stored in $J^\alpha(\tilde w)$ (and similarly for conjecture completions).

Then, we may now define outcome spaces (per completing job):
\[Y^{\alpha}_{\mathrm{proof}} := \{(0,\bot)\}\cup\{(1,d): d\in\mathcal P_{\mathrm{fin}}(\mathcal F)\}\]
\[Y^{\alpha}_{\mathrm{conj}} := \{(0,\bot)\}\cup\{(1,\psi): \psi\in\mathcal F\}\]
And the joint outcome space, dependent on $\tilde w$:
\[Y(\tilde w)
:=
\prod_{\alpha\in\mathsf{Comp}_{\mathrm{proof}}(\tilde w)} Y^{\alpha}_{\mathrm{proof}}
\;\times\;
\prod_{\alpha\in\mathsf{Comp}_{\mathrm{conj}}(\tilde w)} Y^{\alpha}_{\mathrm{conj}}\]

Thus, we may define a product measure over this joint outcome space by conditional independence across completing jobs:
\begin{align*}
&\mathbb K_{\tilde w}(dy)
:=
\\&\left(\prod_{\alpha\in\mathsf{Comp}_{\mathrm{proof}}(\tilde w)}
\mathcal K^\alpha_{\mathrm{proof}}(dy^\alpha \mid \mathcal L^\alpha(\tilde w), s^\alpha_{\mathrm{proof}}(\tilde w),\tau^\alpha_{\mathrm{proof}}(\tilde w))
\right)
\\&\cdot\left(\prod_{\alpha\in\mathsf{Comp}_{\mathrm{conj}}(\tilde w)}
\mathcal K^\alpha_{\mathrm{conj}}(dy^\alpha \mid \mathcal L^\alpha(\tilde w), s^\alpha_{\mathrm{conj}}(\tilde w),\tau^\alpha_{\mathrm{conj}}(\tilde w))
\right)
\end{align*}

\subsubsection{Background - Deterministic Post-Update} 

We now aim to deterministically apply the realized outcomes $y\in Y(\tilde w)$ to update libraries, dependency maps, solve-times, query models, and clear completed jobs.

Given $\tilde w$ and $y\in Y(\tilde w)$, define
\[w^+ := \mathrm{Post}_{\mathrm{nm}}(\tilde w,y)\]
by applying:
\begin{enumerate}
    \item Resolve completing proof jobs. Namely, for each $\alpha\in\mathsf{Comp}_{\mathrm{proof}}(\tilde w)$, let
\[y^\alpha = (b^\alpha, z^\alpha)\in Y^\alpha_{\mathrm{proof}}\]
Let $(s^\alpha,\tau^\alpha)$ be the job’s stored target and total budget. Regardless of proof success, set $U^\alpha_{\mathrm{proof}}(s^\alpha)\ \leftarrow\ \tau^\alpha$.

Then, if $b^\alpha=1$ and $z^\alpha=d$, then update:
\begin{itemize}
    \item	$\mathcal{RS}^\alpha \leftarrow \mathcal{RS}^\alpha \cup \{s^\alpha\}$
	\item	$\delta^\alpha(s^\alpha) \leftarrow d$ (after enforcing $d\subseteq \Gamma$ / acyclicity checks)
	\item	$\Theta^\alpha(s^\alpha)\leftarrow (\tau^\alpha,\alpha)$
	\item	$Q^\alpha \leftarrow \Delta^\alpha(Q^\alpha,(s^\alpha,\Theta^\alpha(s^\alpha)))$
\end{itemize}
	\item Resolve completing conjecture jobs. Namely, for each $\alpha\in\mathsf{Comp}_{\mathrm{conj}}(\tilde w)$, let
\[y^\alpha = (b^\alpha, z^\alpha)\in Y^\alpha_{\mathrm{conj}}\]
Regardless of conjecture success, set $U^\alpha_{\mathrm{conj}}(s^\alpha)\ \leftarrow\tau^\alpha$. Then, if $b^\alpha=1$ and $z^\alpha=\psi$, then:
\[\mathcal C^\alpha \leftarrow \mathcal C^\alpha \cup \{\psi,\neg\psi\}\]
\item Clear completed jobs. Namely, for each $\alpha\in \mathsf{Comp}_{\mathrm{proof}}(\tilde w)\cup \mathsf{Comp}_{\mathrm{conj}}(\tilde w)$, set:
$J^\alpha \leftarrow \bot$.

\end{enumerate}



\subsubsection{Stage I: Non-market transition $\mathcal{T}_{\mathrm{nm}}$}

The non-market transition updates all components of the world state except the market state $\mathcal{M}_t$.
Formally, we define a kernel
\[
\mathcal{T}_{\mathrm{nm}}:\ (\mathcal{W}\times \mathcal{A}_{\mathrm{nm}})\times \Sigma_{\mathcal{W}}\to[0,1]
\]
which is implemented as the following pushforward kernel:
\[\mathcal T_{\mathrm{nm}}(U\mid w,a_{\mathrm{nm}})
\;=\;
\int_{Y(\widetilde w)}
\mathbf 1\!\left[\mathrm{Post}_{\mathrm{nm}}(\widetilde w,y)\in U\right]\,
\mathbb K_{\widetilde w}(dy)\]
For measurable $U\subseteq \mathcal W$, $\widetilde w:=\mathrm{Pre}(w,a_{\mathrm{nm}})$.


\subsubsection{Stage II: Market transition via $\mathscr{M}$}

The market mechanism $\mathscr{M}$ supplies the market transition kernel $\mathcal{T}_{\mathrm{mkt}}$.
Given the realized non-market update $w'_t \sim \mathcal{T}_\mathrm{nm}(\cdot \mid w_t, a_{t,\mathrm{nm}})$ and the market joint action $a_{t,\mathrm{mkt}}$,
we sample the final next state by
\[
w_{t+1}\sim \mathcal{T}_{\mathrm{mkt}}(\cdot\mid w_t,w'_t,a_{t,\mathrm{mkt}}).
\]
Intuitively, $\mathcal{T}_{\mathrm{mkt}}$ may update portfolios $\{\mathcal{P}^\alpha\}$ and the public ledger $\mathcal{ML}$
as a function of both the submitted market actions and the realized mathematical progress encoded in the state change $w_t\to w'_t$,
e.g.\ settlement on newly published proofs, bounty claims, liquidations, AMM price updates, etc. Of course, however, the actual implementation of $\mathcal{T}_\mathrm{mkt}$ is entirely dependent on the specification of the market mechanism $\mathscr{M}$.

\subsubsection{Full transition kernel}


Finally, we may define, for any measurable $U\subseteq\mathcal{W}$,
\begin{align*}
\mathcal{T}(U\mid w_t,a_t) = \int_{\mathcal{W}}
\mathcal{T}_{\mathrm{mkt}}(U\mid w_t,w',a_{t,\mathrm{mkt}})
\ \mathcal{T}_{\mathrm{nm}}(dw'\mid w_t,a_{t,\mathrm{nm}})
\end{align*}



\subsection{Reward, Discount Factor, and Initial Distribution}

We take the per-step reward to be the (expected) \emph{change in balance}:
\[
r^\alpha(w,a) := \mathbb{E}[B^\alpha(w') - B^\alpha(w) \mid w' \sim \mathcal T(\cdot \mid w,a)]
\]

Where $B^\alpha : \mathcal{W} \to \mathbb{R}$ is given by $B^\alpha(w) = \beta(P^\alpha(w))$, where $P^\alpha(w)$ is the usual projection of agent $\alpha$'s portfolio from the $\mathcal{M}$ component of $w \in \mathcal{W}$ and $\beta$ is defined by the $\mathscr{M}$ component of $w \in \mathcal{W}$.

% \subsubsection{Initial Distribution}

Finally, $b_0$ is a distribution over initial world states $w_0\sim b_0$ specifying initial private/shared libraries, initial job statuses
(typically idle), initial solve-time maps (typically empty), initial query models $Q^\alpha_0$, and the initial market state, supported on
states satisfying all invariants.

\subsection{Consistency and well-definedness}

We note that we require a VAMP to be \textit{consistent} with its market mechanism $\mathscr{M}$, meaning that if we have a market mechanism:
\begin{align*}
\mathscr{M} =\Big(&\mathcal{N}, \mathfrak{M},\beta, \mathcal{W},\{\mathcal{A}^\alpha_{\mathrm{mkt}}\}_{\alpha\in\mathcal{N}},
\\&\{\mathrm{Adm}^\alpha_{\mathrm{mkt}}\}_{\alpha\in\mathcal{N}},\mathcal{T}_{\mathrm{mkt}}\Big)
\end{align*}

The following tuple:
\[
\Big(\mathcal{N}',\mathcal{W}',\{\mathcal{O}^\alpha\}_{\alpha\in\mathcal{N}},\{\mathcal{Z}^\alpha\}_{\alpha\in\mathcal{N}},\{\mathcal{A}^\alpha\}_{\alpha\in\mathcal{N}},
\]\[\{\mathrm{Adm}^\alpha\}_{\alpha\in\mathcal{N}},\mathcal{T},
\{r^\alpha\}_{\alpha\in\mathcal{N}},\gamma,b_0, \mathscr{M}\Big)
\]
is a VAMP if and only if $\mathcal N' = \mathcal N, \mathcal W' = \mathcal W$, $\mathcal{A}^\alpha_\mathrm{mkt}\subseteq\mathcal{A}^\alpha$, for $a \in \mathcal{A}^\alpha_\mathrm{mkt}$ and $w \in \mathcal W$,$\mathrm{Adm}^\alpha_\mathrm{mkt}(w,a)=\mathrm{Adm}^\alpha(w,a)$. 

In other words, we may consider the following generator:
\[
\Big(
\mathcal{N},\mathfrak{P},\mathfrak{L}, \beta, \{\mathcal{A}^\alpha_{\mathrm{mkt}}\}_{\alpha\in\mathcal{N}},\{\mathrm{Adm}^\alpha_{\mathrm{mkt}}\}_{\alpha\in\mathcal{N}},\mathcal{T}_{\mathrm{mkt}},\gamma,b_0
\Big)
\]
Which uniquely specifies a well-defined VAMP using the above formulation.




\section{Implementation}
\label{sec:impl}

This section describes a concrete implementation of a restricted subset of VAMP's, with the goal of preserving the core dynamics while making simulation and learning tractable.

\subsection{Assumptions}

We implement a family of environments under the following simplifying assumptions.

\paragraph{Finite Formula Universe.}
Each environment instance is defined over a finite universe of formulas (nodes) $\mathcal{F}$.
All sets derived from $\mathcal{F}$ (ghost/concrete partitions, candidate actions, etc.)
are therefore finite.

\paragraph{Finite Agent Universe.}
We assume $\mathcal{N} = \{1,\cdots,n\}$ for some fixed $n \in \mathbb N$.

% \paragraph{Agent Capabilities.}
% We assume agents $\alpha \in \mathcal{N}$ are equipped with stationary proving and conjecturing abilities. Namely, although agents may learn strategies and update their internal query models based on experience, we freeze their proving and conjecturing kernels: $\mathcal{K}_\mathrm{proof}^\alpha$ $\mathcal{N} = \{1,\cdots,n\}$ for some fixed $n \in \mathbb N$.

\paragraph{Truth and Difficulty Initialization.}

We assume we may initialize the VAMP with a latent truth labeling over formulas $T:\mathcal F\to\{0,1\}$, such that for all $\phi\in\mathcal F$, $T(\phi)+T(\neg\phi)=1$.

We interpret $T(\phi)=1$ as ``$\phi$ is true/valid in the instance'' and $T(\phi)=0$ as ``$\phi$ is false''. For concision, we define $\mathcal{F}_\top := \{\phi \in \mathcal F : T(\phi)=1\}$. Of course, if $[\Gamma \proves \phi]\in\mathcal S$ is a resolved sequent then $\phi \in \mathcal{F}_\top$.
This truth labeling is fixed over time and unobserved by agents except through interaction outcomes.

Additionally, we initialize a static latent difficulty map
\[
d:\mathcal F_\top\to (0,1)
\]
used by the proof kernel as a (unobservable) hyperparameter in assigning success/failure probabilities heterogeneously. $d$ is supported over $\{\phi \in \mathcal F : T(\phi)=1\}$.




\paragraph{Static dependency structure and metadata.}
We assume there exists a fixed, instance-dependent \emph{ground truth} dependency mapping over formulas
\[
\delta^\star : \mathcal{F}_\top\to \mathcal{P}_{\mathrm{fin}}(\mathcal{F})
\]

where we overload this notation to note for resolved sequents $[\Gamma \proves \phi]$ via $\delta^\star([\Gamma \proves \phi]) := \delta^\star(\phi) \subseteq \Gamma$ and such that the induced dependency digraph is acyclic. Then we may define for $s \in \mathcal{RS}_{t_0}^\alpha$ for all $t \ge t_0$, $\delta^\alpha_t(s) = \delta^\star(s)$.

We note this is simply an environment modeling choice. We do not claim mathematical proofs are unique; rather, it yields clean initialization and unambiguous progress signals.

\paragraph{Utility Graph.}
Each instance of a VAMP initializes from a weighted directed graph
\[
G_U := (\mathcal F_\top, E_U, w),\qquad w:E_U\to(0,1],
\]
where $w(\psi,\phi)$ measures the ``utility'' of $\psi$ toward resolving or reasoning about $\phi$. This is necessarily an informal notion that simply serves to allow us to highly related concepts that are not formally dependent on one another. We extend $w$ by $w(\psi,\phi)=0$ when $(\psi, \phi)\notin E_U$.

We additionally note that:
\[w(\psi, \phi)=1 \iff \psi \in \delta^*(\phi)\]
(So, if $E^\star = \{(\psi,\phi) \in E_U : w(\psi,\phi)=1\} \subseteq E_U$, then the subgraph $G^\star = (\mathcal{F}_\top,E^\star)$ is a DAG).


\paragraph{Sequent Specification.}
To prevent an exponential blowup of action space size due to agents' selection of sequents $[\Gamma \proves \phi]$ for $\Gamma \subseteq_{\mathrm{fin}} \mathcal{F}$, we automatically construct antecedent context based on the local library state. Namely, our implemented interface restricts agents to selecting only a consequent formula $\phi$. The antecedent used to evaluate the attempt is then deterministically constructed from the agent's current local library state.

For agent $\alpha$ at time $t$, define the set of locally available resolved formulas
\[
\mathcal{R}^\alpha_t
:= \Big\{\psi\in\mathcal{F}:\exists \Gamma'\subseteq_{\mathrm{fin}}\mathcal{F}\text{ such that }[\Gamma'\proves \psi]\in\mathcal{RS}^\alpha_t\Big\},
\]

Now, we may redefine the (non-market) action space:
\begin{align*}
\widetilde{\mathcal{A}_{\mathrm{nm}}^\alpha}
:=
&
\{\mathrm{Prove}\}\times(\mathcal{F}\times\mathbb{N}^+)
\\&\sqcup
\{\mathrm{Conj}\}\times(\mathcal{F}\times\mathbb{N}^+)
\\&\sqcup
\{\mathrm{Pub}\}\times\mathcal{F}
\\&\sqcup
\{\mathrm{Qry}\}\times \mathcal{S}
\\&\sqcup\{\bot_{\mathrm{nm}}\}
\end{align*}

We note that the query action still relies on sequents rather than formulas, though we will shortly bound the admissible input sequents for this action in particular to allow for polynomial growth of the admissible action space size, while maintaining the ability of agents to ``suppose lemma $\psi$ is proven'' when running a query. W fix a hyperparameter $h_{\mathrm{max}}$ to denote this boundedness of the admissible antecedents' extension beyond $\mathcal{R}^\alpha_t$.

Additionally, as the market mechanism action space is abstractly defined, we note that it is up to the user specification to decide whether similar implementation changes must be changed in $\mathcal{A}_{\mathrm{mkt}}$ as well.

Additionally, we update the specification of the admissibility criteria, where for $w_t \in \mathcal{W}$ and $\alpha \in \mathcal{N}$, if $J^\alpha_t=\bot$, then:
\begin{itemize}
\item For $a^\alpha_t = (\mathrm{Prove},(\phi,\tau))$ and $a^\alpha_t = (\mathrm{Conj},(\phi,\tau))$:
\[
\widetilde{\mathrm{Adm}^\alpha}(w_t,a^\alpha_t)=1 \iff \phi \in \mathcal{C}_t^\alpha
\]
\item For $a^\alpha_t = (\mathrm{Pub},\phi)$: \[\widetilde{\mathrm{Adm}^\alpha}(w_t,a^\alpha_t)=1 \iff \phi \in \mathcal{C}_t^\alpha\]
\item For $a^\alpha_t = (\mathrm{Qry},[\Gamma \proves \phi])$: 
\begin{align*}
    \widetilde{\mathrm{Adm}^\alpha}(w_t,a_t^\alpha)=1 \iff &[\Gamma \proves \phi] \in \mathcal{CS}_t^\alpha \\&\land (|\Gamma \setminus \mathcal{R}^\alpha_t| \le h_{\mathrm{max}})
\end{align*}
\end{itemize}
Where in all other cases, $\widetilde{\mathrm{Adm}^\alpha}(w_t,a_t^\alpha) =\mathrm{Adm}^\alpha(w_t,a_t^\alpha)$. 

Then, given $w_t \in \mathcal{W}$ and admissible $\widetilde{a_t^\alpha} \in \widetilde{\mathcal{A}^\alpha_t}(w_t)$, we may recover a corresponding $a_t^\alpha \in \mathcal{A}^\alpha_t(w_t)$, defined as follows:
\begin{itemize}
    \item If $\widetilde{a_t^\alpha} = (\mathrm{Prove},\phi,\tau)$ then $a_t^\alpha = (\mathrm{Prove},[\mathcal{R}_t^\alpha \proves \phi],\tau)$.
    \item If $\widetilde{a_t^\alpha} = (\mathrm{Conj},\phi,\tau)$ then $a_t^\alpha = (\mathrm{Conj},[\mathcal{R}_t^\alpha \proves \phi],\tau)$.
    \item If $\widetilde{a_t^\alpha} = (\mathrm{Pub},\phi)$ then $a_t^\alpha = (\mathrm{Pub},\{\phi\},\{[\mathcal{R}_t^\alpha \proves \phi]\}\cap \mathcal{R}^\alpha_t)$.
    \item Otherwise, $\widetilde{a_t^\alpha} = a_t^\alpha$.
\end{itemize}


We note that this notion of fixing a canonical local context $\mathcal{R}_t^\alpha$ effectively defines a bijection from (concrete) formulas to (concrete) sequents, in that if we define
\[
s^\alpha_t(\phi)\ :=\ [\Gamma^\alpha_t \proves \phi]\ \in\ \mathcal{S}.
\]
\[
\widetilde{\mathcal{CS}}^\alpha_t
:=\{\,s^\alpha_t(\phi): \phi\in\mathcal{C}^\alpha_t\,\},
\]
and the map $\phi\mapsto s^\alpha_t(\phi)$ is a bijection between $\mathcal{C}^\alpha_t$ and
$\widetilde{\mathcal{CS}}^\alpha_t$. %\textbf{TODO FORMALLY PROVE, AS THIS CS IS DIFFERENT THAN USUAL CS}

Additionally, we note that this assumption automatically implies that all resolved sequents are actually fully resolved.%: \textbf{TODO PROVE!!!}

% We additionally implement the proof kernel with a \emph{resolved-dependency gate}:
% a proof attempt on $\phi$ may only succeed if its returned dependency set $d$ satisfies
% \[
% d\subseteq \Gamma^\alpha_t=\mathrm{ResForm}^\alpha_t.
% \]
% Consequently, any newly resolved sequent $[\Gamma^\alpha_t\proves\phi]$ depends only on formulas that were
% already resolved at time $t$, and therefore (by induction on resolution time, or by acyclicity/topological
% ordering of the dependency DAG) every resolved sequent is automatically \emph{fully} resolved.
% Hence, in the implemented regime we may identify
% \[
% \mathcal{RS}^\alpha_t=\mathcal{FRS}^\alpha_t
% \qquad\text{and similarly for the public library } \mathcal{RS}^0_t=\mathcal{FRS}^0_t.
% \]


\subsection{Proof Kernel ($\mathcal{K}_{\text{proof}}$).}
The proof kernel models the outcome of spending compute budget on an admissible proof attempt. For an agent $\alpha$ and its library state $\mathcal{L}\in\mathscr{L}$, target sequent $s\in\mathcal S$ (though in practice $s \in \mathcal{CS}^{\alpha}$), and compute budget $\tau\in\mathbb N^+$, the kernel outputs a binary success indicator and set of dependencies. However, as we assume a global dependency set, we may simply consider $\mathcal{K}_{\text{proof}}$ outputting in $\{0,1\}$, with the dependency component being computed deterministically as $\delta^\star(s)$.

With that in mind, we now proceed to explicitly define our particular implementation of $\mathcal{K}_{\text{proof}}$, stepping through its parameterization and desired behaviors. Namely, we construct $\mathcal{K}_{\text{proof}}$ to have the following properties:

\paragraph{Bernoulli proof outcomes.}
We assume proof attempts have only binary stochasticity. Concretely, $\mathcal K^\alpha_{\mathrm{proof}}(\cdot\mid \mathcal{L}, s,\tau)$ is supported on exactly two atoms:
\[
\mathcal K^\alpha_{\mathrm{proof}}\big((0,\bot)\mid \mathcal{L}, s,\tau\big)=1-p^\alpha_{\mathrm{proof}}(s,\tau)\]
\[\mathcal K^\alpha_{\mathrm{proof}}\big((1,\delta^\star(s))\mid \mathcal{L}, s,\tau\big)=p^\alpha_{\mathrm{proof}}(s,\tau)
\]
and assigns zero mass to all $(1,d)$ with $d\neq \delta^\star(s)$.

\paragraph{Truth + Dependency Gate.} If $\phi$ is false (i.e. $T(\phi)=0$), the probability of success is $0$. Moreover, if $\delta^\star(\phi)\not\subseteq \mathcal{R}_t^\alpha$, success probability is 0. In other words, one cannot prove a theorem before its hard prerequisites.

Towards this, we may define a feasibility gate
\[
g([\Gamma\proves \phi])=\mathbf{1}[T(\phi)=1]\cdot \mathbf{1}[\delta^\star(\phi)\subseteq \Gamma],
\]

\paragraph{Budget Monotonicity and Diminishing Returns.} We enforce that $p^\alpha_{\mathrm{proof}}(s,\cdot)$ is nondecreasing, intuitively modeling that more cumulative compute can’t hurt success likelihood. Additionally, we hope to model diminishing returns with increasing budget such that marginal gains decrease with more budget. This serves to prevent undesirable strategies such as ``always dump infinite compute into one node''. A natural implementation of this is to model:
\[
p_{\text{proof}}^\alpha([\Gamma \proves \phi],\tau)=K_0\left(1-e^{-K_1\cdot h(\tau)}\right)
\]
For some $K_0,K_1\ge 0$, and a concave $h$. In practice, one often chooses $h(\tau) = \tau^\kappa$ governed by a hyperparameter $\kappa \in (0,1]$.

\paragraph{Difficulty Modulation.} We recall that we initialize $\mathcal{F}_\top$ with a latent difficulty map $d: \mathcal{F}_\top \to (0,1)$. Given the same time budget $\tau$, harder theorems are less likely to be solved. Moreover, difficulty should scale the rate (hazard) rather than just shifting probabilities ad hoc. Towards this notion, define:
\[
r_{\mathrm{diff}}(\phi)=e^{-\lambda_{\mathrm{diff}}\,d(\phi)}
\]
Governed by a difficulty-to-rate transform $\lambda_{\mathrm{diff}}>0$.

\paragraph{Utility Modulation.} We also recall that we initialize a utility graph $G_U$ over $\mathcal{F}_\top$ as a weighted digraph. Although the feasibility gate $g$ filters sequents $[\Gamma \proves \phi]$ based on the inclusion of the hard dependencies $\delta^\star$, we aim to model the notion that even once hard dependencies are met, having additional relevant resolved lemmas in $\Gamma$ should increase success probabilities. This too should be monotone and saturating. Towards this, we define:
\[
s_{\text{mass}}([\Gamma \proves \phi])=\sum_{\psi\in \Gamma} w(\psi,\phi)^\alpha
\]
with $\alpha\ge 1$, and a saturating transform
\[
s([\Gamma \proves \phi])=\log\big(1+s_{\text{mass}}([\Gamma \proves \phi])\big)
\]
This design preserves the distinction between having one useful theorem and many useful theorems, while preventing unbalanced growth in success rate. It is governed by a hyperparameter $\alpha \ge 1$ to affects the sensitivity to high-utility neighbors.

\paragraph{Heterogeneous Agent Capability Modulation.} We initialize agents with capability parameters:
\begin{itemize}
    \item $\rho_0^\alpha\ge 0$ representing agent $\alpha$'s baseline solve-rate
	\item $\rho_1^\alpha\ge 0$ representing agent $\alpha$'s amplification from high-utility support
\end{itemize}
With this in mind, a natural model of the effective proof rate of agent $\alpha$ on a particular target sequent is:
\[
\rho^\alpha([\Gamma \proves \phi])=r_\mathrm{diff}(\phi)\Big(\rho_0^\alpha + \rho_1^\alpha s([\Gamma \proves \phi])\Big)
\]

\subsubsection{Definition}

As such, we may explicitly implement:
\[
\mathcal K^\alpha_{\mathrm{proof}}\big((0,\bot)\mid \mathcal{L}, s,\tau\big)=1-p^\alpha_{\mathrm{proof}}(s,\tau)\]
\[\mathcal K^\alpha_{\mathrm{proof}}\big((1,\delta^\star(s))\mid \mathcal{L}, s,\tau\big)=p^\alpha_{\mathrm{proof}}(s,\tau)
\]
and assigns zero mass to all $(1,d)$ with $d\neq \delta^\star(s)$, where:
\[
p_{\text{proof}}^\alpha([\Gamma \proves \phi],\tau)=g([\Gamma \proves \phi])\left(1-e^{-\rho^\alpha([\Gamma \proves \phi])\cdot h(\tau)}\right)
\]
We note that $p_{\text{proof}}^\alpha(\cdot)\in[0,1]$ for all inputs, and success probability is strictly increasing and approaches $1$ as $\tau_{\mathrm{eff}}\to\infty$ on feasible inputs that pass $g$.

\subsection{Conjecture Kernel ($\mathcal{K}_{\text{conj}}$).}
The conjecture kernel models the outcome of allocating compute budget toward proposing a new formula. Given an agent $\alpha\in\mathcal N$, a library state $\mathcal L=(\mathcal C,\mathcal{RS},\delta,\Theta)\in\mathscr L$, an anchor sequent $s=[\Gamma\proves\phi]\in\mathcal S$ (though in practice $s\in\mathcal{CS}^\alpha$), and budget $\tau\in\mathbb N^+$, the conjecture kernel outputs either failure or success together with a proposed formula $\psi\in\mathcal F$. However, successful outputs are required to land in the current ghost set $\mathcal G(\mathcal{L}) :=\mathcal F\setminus \mathcal C$.

We now define our particular implementation of $\mathcal K^\alpha_{\mathrm{conj}}$, stepping through its desired behaviors and parameterization.

\paragraph{Bernoulli success with conditional proposal.}
We assume conjecturing has binary stochasticity at the success/failure level, followed by a proposal distribution on success. Concretely, there exists a success probability
\[
p^\alpha_{\mathrm{conj}}(\mathcal L,s,\tau)\in[0,1]
\]
and a conditional proposal distribution
\[
\pi^\alpha_{\mathrm{conj}}(\cdot\mid \mathcal L,s,\tau)
\]
supported on $\mathcal G(\mathcal L)$ such that
\[
\mathcal K^\alpha_{\mathrm{conj}}\big((0,\bot)\mid \mathcal L,s,\tau\big)
=
1-p^\alpha_{\mathrm{conj}}(\mathcal L,s,\tau),
\]
and for $\psi\in\mathcal F$,
\[
\mathcal K^\alpha_{\mathrm{conj}}\big((1,\psi)\mid \mathcal L,s,\tau\big)
=
p^\alpha_{\mathrm{conj}}(\mathcal L,s,\tau)\,\pi^\alpha_{\mathrm{conj}}(\psi\mid \mathcal L,s,\tau).
\]

\paragraph{Budget monotonicity and diminishing returns.}
As with the proof kernel, we require $p^\alpha_{\mathrm{conj}}(\mathcal L,s,\cdot)$ to be nondecreasing with diminishing returns. We therefore model
\[
p^\alpha_{\mathrm{conj}}(\mathcal L,s,\tau)
=
1-\exp\!\big(-\eta^\alpha(\mathcal L,s)\,h(\tau)\big),
\]
where $h:\mathbb N^+\to\mathbb R_{\ge 0}$ is increasing and concave. In practice, we again use
\[
h(\tau)=\tau^\kappa,\qquad \kappa\in(0,1].
\]

\paragraph{Heterogeneous agent conjecturing capability.}
We equip each agent $\alpha$ with conjecturing capability parameters
\[
\eta_0^\alpha,\eta_1^\alpha\ge 0,
\]
representing baseline conjecture rate and amplification from local graph opportunity, respectively.

\paragraph{Utility-guided relevance.}
We recall the instance utility graph
\[
G_U=(\mathcal F_\top,E_U,w),
\]
where $w(\psi,\phi)$ measures the utility of $\psi$ toward resolving or reasoning about $\phi$, extended by $0$ off $E_U$.

Intuitively, useful conjectures should depend on the role of the anchor consequent $\phi$ in the current local library. If $\phi$ is already resolved, then good conjectures should tend to extend outward from $\phi$; if $\phi$ is not yet resolved, then good conjectures should tend to provide lemmas that support $\phi$.

% Accordingly, define the local resolved formula set induced by $\mathcal{RS}$:
% \[
% \mathcal R(\mathcal L):=\Big\{\psi\in\mathcal F:\exists \Gamma'\subseteq_{\mathrm{fin}}\mathcal F,\; [\Gamma'\proves\psi]\in\mathcal{RS}\Big\},
% \]
and define the anchor-conditioned utility score
\[
q(\psi\mid \mathcal L,[\Gamma\proves\phi])
=
\mathbf 1[[\Gamma\proves\phi]\in\mathcal {RS}^\alpha]\,w(\phi,\psi)
+
\mathbf 1[[\Gamma\proves\phi]\notin\mathcal {RS}^\alpha]\,w(\psi,\phi).
\]
Thus, when $\phi$ is resolved, the score favors frontier extensions; when $\phi$ is unresolved, it favors potentially useful supporting lemmas.

\paragraph{Local opportunity modulation.}
We also want conjecturing to be easier in regions where there exist many promising unseen formulas. To capture this, define the local opportunity mass
\[
m^\alpha(\mathcal L,s)
:=
\sum_{\psi\in\mathcal G(\mathcal L)} \big(q(\psi\mid \mathcal L,s)\big)^{\beta_{\mathrm{conj}}},
\qquad
\beta_{\mathrm{conj}}\ge 1,
\]
and the saturating transform
\[
u^\alpha(\mathcal L,s):=\log\big(1+m^\alpha(\mathcal L,s)\big).
\]
We then define the effective conjecturing rate by
\[
\eta^\alpha(\mathcal L,s)
=
\eta_0^\alpha+\eta_1^\alpha\,u^\alpha(\mathcal L,s).
\]

\paragraph{Budget-improved proposal quality.}
Beyond increasing the chance of any successful conjecture, additional compute should also improve the quality of the proposed conjecture conditional on success. We model this by allowing the proposal distribution to become more concentrated around high-utility candidates as budget increases.

To do so, define a budget-dependent selectivity parameter
\[
\kappa^\alpha_{\mathrm{conj}}(\tau)
=
\kappa^\alpha_0+\kappa^\alpha_1 h(\tau),
\qquad
\kappa^\alpha_0,\kappa^\alpha_1\ge 0.
\]
Small values of $\kappa^\alpha_{\mathrm{conj}}(\tau)$ yield diffuse proposal distributions, while large values concentrate mass on high-score candidates.

\paragraph{Proposal distribution over all formulas, restricted to ghosts by renormalization.}
It is convenient to first define a raw proposal score over all formulas $\psi\in\mathcal F$, then restrict to the current ghost set induced by $\mathcal L$.

Let
\[
r^\alpha(\psi\mid \mathcal L,s,\tau)
:=
\exp\!\Big(\kappa^\alpha_{\mathrm{conj}}(\tau)\,\varphi(q(\psi\mid \mathcal L,s))\Big),
\]
where $\varphi:[0,1]\to\mathbb R_{\ge0}$ is a fixed monotone transform (for instance $\varphi(x)=x$ or $\varphi(x)=\log(1+x)$).

We then define the actual proposal distribution by renormalizing over the current ghost set:
\[
\pi^\alpha_{\mathrm{conj}}(\psi\mid \mathcal L,s,\tau)
=
\frac{r^\alpha(\psi\mid \mathcal L,s,\tau)\,\mathbf 1[\psi\in\mathcal G(\mathcal L)]}
{\sum_{\psi'\in\mathcal G(\mathcal L)} r^\alpha(\psi'\mid \mathcal L,s,\tau)}.
\]
Hence the conjecturer is specified globally over $\mathcal F$, but at runtime only samples from formulas that are currently ghost relative to $\mathcal L$.

\subsubsection{Definition}

Putting the above together, for an admissible conjecture attempt on $s=[\Gamma\proves\phi]$ with budget $\tau$, we define
% \[
% p^\alpha_{\mathrm{conj}}(\mathcal L,s,\tau)
% =
% 1-\exp\!\big(-\eta^\alpha(\mathcal L,s)\,h(\tau)\big),
% \]
% with
% \[
% \eta^\alpha(\mathcal L,s)
% =
% \eta_0^\alpha+\eta_1^\alpha\,\log\!\Big(1+\sum_{\psi\in\mathcal G(\mathcal L)}(q(\psi\mid \mathcal L,s))^{\beta_{\mathrm{conj}}}\Big),
% \]
% and
% \[
% \pi^\alpha_{\mathrm{conj}}(\psi\mid \mathcal L,s,\tau)
% =
% \frac{\exp\!\Big(\kappa^\alpha_{\mathrm{conj}}(\tau)\,\varphi(q(\psi\mid \mathcal L,s))\Big)\,\mathbf 1[\psi\in\mathcal G(\mathcal L)]}
% {\sum_{\psi'\in\mathcal G(\mathcal L)}
% \exp\!\Big(\kappa^\alpha_{\mathrm{conj}}(\tau)\,\varphi(q(\psi'\mid \mathcal L,s))\Big)}.
% \]

% Accordingly,
\[
\mathcal K^\alpha_{\mathrm{conj}}\big((0,\bot)\mid \mathcal L,s,\tau\big)
=
1-p^\alpha_{\mathrm{conj}}(\mathcal L,s,\tau),
\]
and for $\psi\in\mathcal F$,
\[
\mathcal K^\alpha_{\mathrm{conj}}\big((1,\psi)\mid \mathcal L,s,\tau\big)
=
p^\alpha_{\mathrm{conj}}(\mathcal L,s,\tau)\,\pi^\alpha_{\mathrm{conj}}(\psi\mid \mathcal L,s,\tau).
\]
By construction, $\pi^\alpha_{\mathrm{conj}}(\cdot\mid \mathcal L,s,\tau)$ is supported on $\mathcal G(\mathcal L)$, $p^\alpha_{\mathrm{conj}}(\mathcal L,s,\cdot)$ is monotone with diminishing returns, and larger budgets increase both the chance of successful conjecturing and the concentration of mass on high-utility conjectures.

\subsection{Query Model.}

The query model is a stateful predictive module used by each agent to estimate how its prover process is expected to perform on a target sequent under its current local library state. In our implementation, the query model does \emph{not} estimate semantic truth directly. Rather, for agent $\alpha$, it estimates:

\begin{enumerate}
    \item the probability that $\alpha$ will \emph{eventually resolve} a target sequent under its prover process, and
    \item the expected \emph{additional compute budget} required to do so, conditional on eventual success.
\end{enumerate}

Accordingly, we instantiate the query readout as a map
\[
q^\alpha : \mathcal Q^\alpha \times \mathscr L \times \mathcal S \to [0,1]\times \mathbb N,
\]
where
\[
q^\alpha(Q^\alpha_t,\mathcal L,s)
=
\big(\hat p^\alpha_t(\mathcal L,s),\hat \tau^\alpha_t(\mathcal L,s)\big).
\]
Here:
\begin{itemize}
    \item $\hat p^\alpha_t(\mathcal L,s)$ estimates the probability that agent $\alpha$'s prover process will eventually resolve $s$ from local library state $\mathcal L$,
    \item $\hat\tau^\alpha_t(\mathcal L,s)$ estimates the expected additional budget required to resolve $s$, conditional on eventual success.
\end{itemize}

\paragraph{Feature buckets.}
To make online updates statistically stable while preserving algebraic simplicity, we use a finite bucketization of query instances. For each agent $\alpha$, fix a finite bucket set $\mathcal B^\alpha$ and a deterministic bucket map
\[
b^\alpha : \mathscr L\times \mathcal S \to \mathcal B^\alpha.
\]
The bucket map may depend on any collection of state/query features, such as:
\begin{itemize}
    \item support mass or utility features of the target,
    \item dependency-completion indicators,
    \item target difficulty proxies,
    \item cumulative prior budget spent on the target,
    \item local graph-topological features,
    \item other agent-specific or state-specific summary statistics.
\end{itemize}
The framework does not prescribe the exact feature design.

\paragraph{Discrete-time survival model.}
Fix a finite horizon $H\in\mathbb N^+$. For each agent $\alpha$ and bucket $b\in\mathcal B^\alpha$, we model a discrete-time hazard sequence
\[
\lambda^\alpha_b(k)\in[0,1], \qquad k=1,\dots,H,
\]
where
\[
\lambda^\alpha_b(k)
:=
\Pr(T^\alpha_s = k \mid T^\alpha_s \ge k,\ b^\alpha(\mathcal L,s)=b).
\]
Here $T^\alpha_s$ denotes the random additional budget required for agent $\alpha$ to resolve $s$.

Equivalently, conditional on bucket $b$, the discrete survival function is
\[
S^\alpha_b(k)
:=
\Pr(T^\alpha_s > k \mid b)
=
\prod_{j=1}^{k}\big(1-\lambda^\alpha_b(j)\big),
\qquad S^\alpha_b(0):=1.
\]
Thus the probability of solving exactly at budget step $k$ is
\[
\Pr(T^\alpha_s = k \mid b)
=
\lambda^\alpha_b(k)\prod_{j=1}^{k-1}(1-\lambda^\alpha_b(j)).
\]

\paragraph{Query readout.}
Given current state $(Q^\alpha_t,\mathcal L,s)$, let
\[
b=b^\alpha(\mathcal L,s).
\]
Then the query model returns:
\[
\hat p^\alpha_t(\mathcal L,s)
=
\Pr(T^\alpha_s \le H \mid b)
=
1-S^\alpha_b(H)
=
1-\prod_{j=1}^{H}(1-\lambda^\alpha_b(j)),
\]
which is the estimated probability of eventual solve success within the modeling horizon $H$.

For the conditional expected additional budget, define
\[
\hat\tau^\alpha_t(\mathcal L,s)
=
\mathbb E[T^\alpha_s \mid T^\alpha_s \le H,\ b]
=
\frac{\sum_{k=1}^{H} k\,\Pr(T^\alpha_s = k \mid b)}
{\hat p^\alpha_t(\mathcal L,s)},
\]
whenever $\hat p^\alpha_t(\mathcal L,s)>0$. If $\hat p^\alpha_t(\mathcal L,s)=0$, we may return a fixed sentinel value (for instance $H$ or $+\infty$ in an abstract implementation).

\paragraph{Query model state.}
For each agent $\alpha$, the query model state
\[
Q^\alpha_t \in \mathcal Q^\alpha
\]
stores, for every bucket $b\in\mathcal B^\alpha$ and time index $k\in\{1,\dots,H\}$, two nonnegative sufficient statistics:
\[
N^\alpha_{t}(b,k),\qquad D^\alpha_{t}(b,k).
\]
These are interpreted as:
\begin{itemize}
    \item $D^\alpha_t(b,k)$ = the number of observed proof-attempt trajectories in bucket $b$ that remained ``at risk'' at budget step $k$,
    \item $N^\alpha_t(b,k)$ = the number of those trajectories that terminated successfully exactly at budget step $k$.
\end{itemize}

\paragraph{MAP-smoothed hazard estimates.}
To ensure numerical stability and prevent degenerate zero/one hazards early in training, we introduce fixed hyperparameters
\[
a^\alpha_{b,k}>0,\qquad c^\alpha_{b,k}>0
\]
for each bucket-time pair. These act as pseudocounts. We then define the hazard estimate by the smoothed ratio
\[
\lambda^\alpha_b(k)
=
\frac{N^\alpha_t(b,k)+a^\alpha_{b,k}}
{D^\alpha_t(b,k)+a^\alpha_{b,k}+c^\alpha_{b,k}}.
\]
This may be interpreted as the coordinatewise MAP estimate of a Bernoulli hazard with Beta prior.

\paragraph{Observed data and censoring.}
Each proof attempt produces an observation datum for the querying agent. We model a datum as
\[
d=(b,\tau,\xi)\in \mathcal B^\alpha \times \mathbb N^+ \times \{0,1\},
\]
where:
\begin{itemize}
    \item $b$ is the feature bucket of the queried instance at attempt time,
    \item $\tau$ is the additional budget consumed in that attempt,
    \item $\xi=1$ indicates that the attempt succeeded exactly at budget $\tau$,
    \item $\xi=0$ indicates that the attempt was censored (timed out / did not solve) at budget $\tau$.
\end{itemize}

Such a datum contributes additive increments to the sufficient statistics:
\[
D^\alpha_t(b,k) \leftarrow D^\alpha_t(b,k)+1
\qquad\text{for all }1\le k\le \min\{\tau,H\},
\]
and, if $\xi=1$ and $\tau\le H$, additionally
\[
N^\alpha_t(b,\tau) \leftarrow N^\alpha_t(b,\tau)+1.
\]
Hence unsuccessful attempts contribute exposure only, while successful attempts contribute both exposure and one event at the realized solve time.

\paragraph{Update rule.}
Define the update operator
\[
\Delta^\alpha : \mathcal Q^\alpha \times \mathcal D(\mathcal Q^\alpha)\to \mathcal Q^\alpha
\]
to be the additive update of these sufficient statistics induced by a single datum $d$ as above. Because the update is coordinatewise addition on count tables, it is exactly associative and commutative. Therefore, if multiple observations are received in a single timestep, the resulting batch update is well-defined independently of order.

\subsection{Market Mechanisms}

We now specify the family of concrete market mechanisms tested in this work. Each mechanism is implemented as an instantiation of the abstract market mechanism interface
\[
\mathscr M=\Big(\mathcal N,\mathfrak M,\beta,\mathcal W,\{\mathcal A^\alpha_{\mathrm{mkt}}\}_{\alpha\in\mathcal N},
\{\mathrm{Adm}^\alpha_{\mathrm{mkt}}\}_{\alpha\in\mathcal N},\mathcal T_{\mathrm{mkt}}\Big),
\]
but we restrict attention to a common implementation template so that different mechanisms may be compared under the same underlying theorem-proving environment.

\paragraph{Common portfolio structure.}
For all tested mechanisms, each agent $\alpha$ maintains a portfolio
\[
\mathcal P^\alpha_t=(c^\alpha_t,H^\alpha_t,O^\alpha_t),
\]
where:
\begin{itemize}
    \item $c^\alpha_t\in\mathbb R_{\ge 0}$ is current cash,
    \item $H^\alpha_t$ is the multiset of currently held assets / claims,
    \item $O^\alpha_t$ is the multiset of currently posted market offers.
\end{itemize}
The public market ledger $\mathcal{ML}_t$ stores the union of all currently active public offers together with any mechanism-specific public state (for example AMM inventory variables or posted bounty books).

\paragraph{Worst-case collateralization.}
To prevent agents from taking positions they cannot honor, we associate to each portfolio a worst-case marked balance
\[
\underline B^\alpha_t := \underline\beta(\mathcal P^\alpha_t)\in\mathbb R,
\]
defined as the cash position minus the largest possible future liability of the held and posted claims under the mechanism's settlement rule. We require all admissible market actions to preserve
\[
\underline B^\alpha_t \ge 0.
\]
Thus agents are fully collateralized against worst-case future settlement.

\paragraph{Cash conservation.}
All tested mechanisms satisfy internal money conservation at trade time: transferring an asset from one agent to another is accompanied by an equal and opposite transfer of cash, except where the mechanism explicitly introduces escrowed collateral or deterministic settlement transfers already implied by outstanding liabilities.

\paragraph{Public settlement events.}
All tested mechanisms settle only against publicly observable events, namely changes in the public library $\mathcal L^0_t$. In particular, contract payouts may depend only on whether a target formula or sequent has been publicly resolved by a specified time.


\paragraph{Resolution events.}
For a public target sequent $s$ and deadline $T\in\mathbb N$, define the binary event
\[
E_{s,T}(w_{0:\infty}) := \mathbf 1\big[\exists\, t\le T \text{ such that } s\in \mathcal{RS}^0_t\big].
\]
Equivalently, $E_{s,T}=1$ iff $s$ is publicly resolved by time $T$.
When working in the implemented formula-based interface, we identify a formula $\phi$ with its canonical public sequent
\[
s^0_t(\phi):=[\mathcal R^0_t\proves \phi]
\]
and write $E_{\phi,T}$ by abuse of notation.




\subsubsection{Mechanism I: Collateralized Bilateral Contract Market}

Our first mechanism is a posted-offer market over binary contingent contracts. A \emph{contract type} is a triple
\[
\chi=(s,T,\ell)\in \mathcal{RS}^{0,\mathrm{cand}}\times \mathbb N \times [0,1],
\]
where:
\begin{itemize}
    \item $s$ is a public target sequent,
    \item $T$ is a strike / deadline time,
    \item $\ell\in[0,1]$ is the maximum loss parameter.
\end{itemize}

Intuitively, $\chi=(s,T,\ell)$ is a normalized unit-notional claim on the event that $s$ becomes publicly resolved by time $T$.

\paragraph{Long and short claims.}
Each contract type $\chi$ induces two complementary positions:
\[
\chi^{+} \qquad\text{and}\qquad \chi^{-},
\]
called the long and short side, respectively. Their terminal cash payoffs are:
\[
\Pi(\chi^+) =
\begin{cases}
1-\ell, & E_{s,T}=1,\\
-\ell, & E_{s,T}=0,
\end{cases}
\qquad
\Pi(\chi^-) = -\Pi(\chi^+).
\]
Thus the long side gains when the target resolves by the deadline, while the short side gains otherwise. The magnitude of the downside of the long position is exactly $\ell$, and the upside is $1-\ell$.

\paragraph{Unit notional and quantities.}
Contracts are traded in nonnegative integer quantities. A position of quantity $q\in\mathbb N$ in $\chi^+$ has payoff $q\,\Pi(\chi^+)$, and similarly for $\chi^-$. Hence larger notional exposure is represented by larger quantities rather than by changing the contract definition itself.

\paragraph{Asset creation.}
A new quantity-$q$ contract issue of type $\chi$ is created only as a matched pair:
\[
q\chi^+ \quad\text{and}\quad q\chi^-.
\]
This pair may initially be placed into the issuing agent's holdings or posted offers. Since the long and short payoffs sum to zero statewise, pair creation introduces no net wealth.

\paragraph{Public offers.}
An offer is a tuple
\[
o=(\chi,\sigma,q,p),
\]
where $\chi$ is a contract type, $\sigma\in\{+,-\}$ is the side being offered, $q\in\mathbb N^+$ is quantity, and $p\in[0,1]$ is the per-unit transaction price. Posted offers are stored in the public ledger $\mathcal{ML}_t$.

\paragraph{Trades.}
When agent $\beta$ accepts quantity $q'\le q$ from an offer $o=(\chi,\sigma,q,p)$ posted by agent $\alpha$, the ledger removes $q'\chi^\sigma$ from $\alpha$'s posted offers and transfers it to $\beta$'s holdings, while cash is transferred in the opposite direction:
\[
c^\beta \leftarrow c^\beta - q'p,\qquad c^\alpha \leftarrow c^\alpha + q'p.
\]
The residual quantity $q-q'$ remains posted if positive.

\paragraph{Collateral and admissibility.}
For a held or posted position $x$, define its worst-case liability by
\[
L(x):=\sup_{\omega}(-\Pi_\omega(x))_+.
\]
For a portfolio $\mathcal P^\alpha=(c^\alpha,H^\alpha,O^\alpha)$, define the minimum marked balance
\[
\underline B^\alpha
:=
c^\alpha - \sum_{x\in H^\alpha\cup O^\alpha}L(x).
\]
A market action is admissible only if the resulting portfolio satisfies
\[
\underline B^\alpha \ge 0.
\]
Hence every agent is fully collateralized against the worst-case settlement of all currently held and posted obligations.

\paragraph{Settlement.}
At each timestep, any contract type $\chi=(s,T,\ell)$ with deadline $T=t$ is settled against the public event $E_{s,T}$. For each agent, the corresponding cash payoff is applied to all held and posted quantities of $\chi^\pm$, and the settled contracts are removed from portfolios and the public ledger.





\subsubsection{Mechanism II: Compositional Contract Market}

Our second mechanism generalizes Mechanism I by allowing agents to trade a richer family of contingent assets built compositionally from primitive public resolution events.

\paragraph{Primitive event contracts.}
For each public target sequent $s$ and deadline $T$, let
\[
X_{s,T}:=E_{s,T}\in\{0,1\}
\]
denote the primitive event indicator that $s$ is publicly resolved by time $T$.

\paragraph{Composable payoff language.}
Fix a class $\mathfrak C$ of admissible payoff expressions generated inductively from:
\begin{itemize}
    \item constants $a\in\mathbb R$,
    \item primitive event indicators $X_{s,T}$,
    \item affine combinations,
    \item pointwise minimum / maximum,
    \item and other chosen bounded payoff constructors.
\end{itemize}
Each contract expression $C\in\mathfrak C$ induces a bounded terminal payoff
\[
\Pi(C)\in\mathbb R
\]
which is a deterministic measurable function of the public resolution path.

\paragraph{Examples.}
This language includes, for example:
\begin{itemize}
    \item binary event contracts,
    \item deadline spreads,
    \item piecewise-linear payout schedules,
    \item capped bundles of multiple theorem-resolution events.
\end{itemize}

\paragraph{Trading and collateralization.}
Trading proceeds exactly as in Mechanism I, except that contract types are replaced by general payoff expressions $C\in\mathfrak C$. Worst-case collateral is computed from the statewise lower bound of the payoff:
\[
L(C):=\sup_\omega(-\Pi_\omega(C))_+.
\]
Admissibility again requires full collateralization under worst-case settlement.

\paragraph{Purpose.}
This mechanism allows us to test whether richer contingent-claim design improves allocation efficiency relative to simple binary contracts, while preserving the same posted-offer and collateral framework.


\subsubsection{Mechanism III: LMSR Prediction Market}

Our third mechanism is an automated market maker over binary public-resolution events using a logarithmic market scoring rule (LMSR).

\paragraph{Binary event markets.}
For each public target event $(s,T)$, define a binary market on
\[
E_{s,T}\in\{0,1\}.
\]
The public ledger stores an inventory state
\[
\mathbf q_{s,T}=(q^1_{s,T},q^0_{s,T})\in\mathbb R^2,
\]
where $q^1_{s,T}$ and $q^0_{s,T}$ are the outstanding share inventories of the YES and NO sides, respectively.

\paragraph{Cost function.}
Fix a liquidity parameter $b>0$. The LMSR cost function for market $(s,T)$ is
\[
C_{s,T}(\mathbf q)
=
b\log\!\Big(\exp(q^1/b)+\exp(q^0/b)\Big).
\]

\paragraph{Trade pricing.}
If an agent changes inventory from $\mathbf q$ to $\mathbf q+\Delta \mathbf q$, then the cash paid to the market maker is
\[
C_{s,T}(\mathbf q+\Delta \mathbf q)-C_{s,T}(\mathbf q).
\]
Thus marginal prices are
\[
p^1_{s,T}(\mathbf q)
=
\frac{\exp(q^1/b)}{\exp(q^1/b)+\exp(q^0/b)},
\qquad
p^0_{s,T}(\mathbf q)=1-p^1_{s,T}(\mathbf q).
\]

\paragraph{Holdings and settlement.}
An agent may hold YES and NO shares in each event market. At settlement time $T$, YES shares pay $1$ iff $E_{s,T}=1$, and NO shares pay $1$ iff $E_{s,T}=0$.

\paragraph{Admissibility.}
An agent trade is admissible only if its post-trade cash remains nonnegative (or, in a collateralized extension, above a specified lower bound). Because the LMSR market maker internalizes the pricing rule, bilateral counterparties are not required.

\paragraph{Purpose.}
This mechanism provides a smooth prediction-market baseline with endogenous prices and continuous liquidity, allowing us to compare decentralized theorem-allocation behavior against the bilateral posted-offer markets above.

\subsubsection{Mechanism IV: Posted Bounty Market}

Our fourth mechanism removes secondary trading entirely and instead allocates rewards through posted bounties on public resolution events.

\paragraph{Bounty objects.}
A bounty is a tuple
\[
b=(s,T,R,\alpha),
\]
where:
\begin{itemize}
    \item $s$ is a public target sequent,
    \item $T$ is an optional deadline,
    \item $R>0$ is the reward amount,
    \item $\alpha$ is the sponsoring agent who escrows the reward.
\end{itemize}

\paragraph{Escrow.}
When agent $\alpha$ posts bounty $b=(s,T,R,\alpha)$, the amount $R$ is removed from liquid cash and placed in escrow. This prevents the sponsor from withdrawing the reward before settlement.

\paragraph{Winner-take-all settlement.}
If $s$ becomes publicly resolved by time $T$, then the bounty is paid to the first agent whose publication caused $s$ to enter the public resolved set. Formally, if
\[
t^\star := \min\{t\le T : s\in\mathcal{RS}^0_t\},
\]
and $\Theta^0_{t^\star}(s)=(\tau,\beta)$ records $\beta$ as the publishing / resolving agent, then $\beta$ receives $R$ at settlement. If no such public resolution occurs by time $T$, then the bounty expires and the escrowed amount is returned to the sponsor, or alternatively burned, depending on the chosen implementation variant.

\paragraph{Bounty ledger.}
The public market ledger stores all active bounties. Agents may condition their proving, conjecturing, and publication strategies on this posted reward landscape.

\paragraph{Purpose.}
This mechanism provides a sparse incentive baseline in which only the eventual public resolver is rewarded, in contrast to the tradable-claims markets where many agents may profit from correct early positioning.

\section{Multiagent Reinforcement Learning}
\subsection{Architecture}

We model each agent as a policy over a partially observable trajectory and train under a centralized-training/decentralized-execution (CTDE) paradigm. In the current prototype, the actor is implemented as a \emph{causal transformer} over a serialized history of structured tokens rather than a graph neural network or set-specific architecture. We choose this design for rapid prototyping on small synthetic instances (up to roughly $10^3$ formula nodes per environment instance), where a standard transformer is computationally feasible and easier to iterate on.

\paragraph{Trajectory Tokenization.}
For each agent $a$, we construct a recent history window of length $H$ timesteps and serialize it as
\[
[\mathrm{OBS}_{t-H+1}^{(a)}][\mathrm{ACT}_{t-H+1}^{(a)}][\mathrm{EVT}_{t-H+1\rightarrow t-H+2}]
\cdots
[\mathrm{OBS}_t^{(a)}],
\]
where:
\begin{itemize}
    \item $\mathrm{OBS}_\tau^{(a)}$ is the agent's local observation at time $\tau$,
    \item $\mathrm{ACT}_\tau^{(a)}$ is the agent's executed action at time $\tau$,
    \item $\mathrm{EVT}_{\tau\rightarrow \tau+1}$ is a (possibly empty) sequence of \emph{public} event tokens that summarize observable actions and market updates by all agents between decision points.
\end{itemize}
This representation distinguishes the agent's own control decisions from exogenous public dynamics while preserving temporal order.

\paragraph{Observation Block Structure.}
Each observation block is serialized as
\begin{align*}
    [\mathrm{OBS\_START}][\mathrm{GLOBAL}]\\
[\mathrm{LIB\_START}]\;\texttt{lib\_item}_1\cdots \texttt{lib\_item}_{m_\tau}\;[\mathrm{LIB\_END}]\\
[\mathrm{PORT\_START}]\;\texttt{balance}\;\texttt{asset}_1\cdots \texttt{asset}_{p_\tau}\;[\mathrm{PORT\_END}]\\
[\mathrm{MKT\_START}]\;\texttt{offer}_1\cdots \texttt{offer}_{q_\tau}\;[\mathrm{MKT\_END}]\\
[\mathrm{OBS\_END}],
\end{align*}
where $m_\tau,p_\tau,q_\tau$ vary across timesteps.

For a concrete theorem node $i$, the library token is
\[
\texttt{lib\_item}(i,\tau)=\mathbf{e}_{\mathrm{id}}(i)+\phi_{\mathrm{node}}(f_i^\tau),
\]
where $\mathbf{e}_{\mathrm{id}}(i)$ is a fixed non-learned PRNG/hash embedding used as an environment-local identifier (a ``nametag''), and $\phi_{\mathrm{node}}$ is an MLP over structured node features
\[
f_i^\tau = (\texttt{isResolved},\texttt{isPublic},\texttt{isNegated},\texttt{ageConj},\texttt{ageRes}).
\]
Here, $\texttt{ageConj}$ and $\texttt{ageRes}$ are discretized or clipped ages relative to the current timestep rather than absolute timestamps.

For a portfolio-held asset or public offer referencing theorem $i$, the token is
\[
\texttt{asset}(i,\tau)=\mathbf{e}_{\mathrm{id}}(i)+\phi_{\mathrm{asset}}(g_i^\tau),
\]
where $g_i^\tau$ includes structured contract and bookkeeping features such as side (buy/sell), time-to-strike, quantity, max-loss parameter, and flags indicating whether the asset is privately held, a public offer, or the agent's own posted offer. A separate balance token encodes scalar accounting state (e.g., cash balance and optional reserved collateral) via a small MLP.

\paragraph{Action and Event Tokens.}
We represent each executed action as a single structured token
\[
\texttt{act}_\tau=\mathbf{e}_{\mathrm{type}}(a_\tau)+\mathbf{e}_{\mathrm{target}}+\phi_{\mathrm{act}}(\texttt{params}),
\]
where $\mathbf{e}_{\mathrm{type}}$ is a learned action-type embedding, $\mathbf{e}_{\mathrm{target}}$ is either the identifier embedding of the targeted theorem/offer (when applicable) or a null vector, and $\phi_{\mathrm{act}}$ embeds numeric action parameters (e.g., proof budget, conjecture budget, publication heuristic parameters, or offer parameters).

Public events (e.g., publications, posted offers, accepted offers, settlements) are tokenized analogously using event-type embeddings plus target identifiers and event parameters. This enables the model to condition on strategic behavior by other agents without conflating those dynamics with the acting agent's own actions.

\paragraph{Actor Network.}
The token sequence is embedded into a shared model dimension and processed by a causal transformer with learned positional encodings. Let $h_t$ denote the final hidden state associated with the last token (or a dedicated prediction token appended after $\mathrm{OBS}_t^{(a)}$). The actor maps $h_t$ to a structured action distribution via multiple heads:
\begin{enumerate}
    \item an \textbf{action-type head} over the discrete action family (conjecture, prove, publish, market actions, etc.),
    \item a \textbf{target-selection head} implemented as a pointer-style distribution over the relevant current candidates (e.g., theorem nodes in the current library block or offers in the current market block),
    \item one or more \textbf{parameter heads} for budgets and contract parameters (e.g., compute allocation, strike-time offset, quantity, max-loss).
\end{enumerate}
This factorization preserves a compact policy output while supporting variable-size candidate sets.

\paragraph{Centralized Critic.}
During training, we use a centralized critic that conditions on a richer global summary than any individual agent can observe. In the current implementation, the critic receives environment-level aggregate features (e.g., statistics of all agents' private/public libraries, market summaries, and graph-generation metadata) along with the joint public state, and outputs value estimates used for policy optimization. At execution time, only the decentralized actor is used.
\subsubsection{Training}

Our training pipeline combines offline initialization with online multiagent reinforcement learning under CTDE. The transformer-based actor is trained as a policy network (not purely as a sequence-modeling planner), while retaining the benefits of a structured trajectory representation for long-horizon credit assignment.

\paragraph{Stage 1: Offline Initialization (Behavior Cloning / Warm Start).}
We initialize actor policies on logged trajectories generated by scripted and heuristic agents. These trajectories are serialized using the same tokenization scheme employed during online training (observation blocks, action tokens, and public event tokens), allowing direct reuse of the actor architecture during warm start.

The scripted baselines include:
\begin{itemize}
    \item dependency-aware or feasibility-first proving heuristics,
    \item utility-greedy conjecturing heuristics,
    \item immediate vs delayed publication heuristics,
    \item simple market-making / bounty-following / claim-arbitrage heuristics (when market modules are enabled).
\end{itemize}
We train the actor to predict the scripted action type and action parameters (including target selections) with supervised losses. This warm start reduces exploration burden and decreases the likelihood of early collapse into degenerate behaviors (e.g., random conjecturing or chronically delayed publication).

\paragraph{Stage 2: Online CTDE Fine-Tuning (MAPPO-style).}
We then fine-tune the warm-started actor with online multiagent RL using PPO-style policy optimization under CTDE (MAPPO-style updates for decentralized actors and a centralized critic). Rollouts are collected from populations of interacting agents in procedurally generated Forge instances.

At each update cycle, we:
\begin{enumerate}
    \item sample a batch of environment instances (with controlled graph/topology parameters),
    \item run multiagent rollouts using the current policy population,
    \item compute returns/advantages using the centralized critic,
    \item update decentralized actor parameters with clipped PPO objectives,
    \item update the centralized critic on global-state value targets.
\end{enumerate}
We treat action legality constraints (e.g., invalid targets, inadmissible proof attempts) as environment-enforced and optionally apply action masking in the policy heads to improve sample efficiency.

\paragraph{Population Training and Opponent Diversity.}
To reduce overfitting to a single co-adapted equilibrium, we train against a population of agents rather than a single self-play copy. In practice, we maintain a replayable pool containing current policies, lagged checkpoints, and heuristic agents, and sample opponents from this pool during rollout generation. This approximates league-style training and is intended to improve robustness under non-stationary multiagent dynamics.

\paragraph{Curriculum Learning.}
Because the full environment induces large candidate sets and long credit-assignment horizons, we use a curriculum over instance complexity. Training begins on small environments with shallow dependency DAGs, fewer agents, and limited market mechanics, then progressively increases:
\begin{itemize}
    \item number of formula nodes,
    \item DAG depth and branching factor,
    \item number of agents,
    \item market complexity (bounties/claims/offers/publication timing),
    \item degree of private information asymmetry.
\end{itemize}
The current prototype architecture (single causal transformer over serialized token histories) is designed for the early and mid stages of this curriculum; larger-scale versions may replace the observation encoder with a set-structured architecture to improve permutation invariance and scalability.
\subsubsection{Evaluation}

We evaluate this system along both \emph{environmental} and \emph{agentic} axes.


\paragraph{Agent-Level Performance.}
For learned agents, we measure:
\begin{itemize}
    \item cumulative utility and regret relative to scripted baselines,
    \item number and value-weighted count of resolved theorems,
    \item proof efficiency (utility per unit compute),
    \item conjecture quality (average utility relevance of generated conjectures),
    \item publication efficiency and private-advantage retention.
\end{itemize}

\paragraph{System-Level Outcomes.}
At the market/game level, we evaluate:
\begin{itemize}
    \item total resolved theorem utility over time,
    \item duplication rate of effort across agents,
    \item time-to-resolution for high-value targets,
    \item credit assignment quality (alignment between causal contribution and reward),
    \item emergent specialization and division of labor.
\end{itemize}

Since this report describes in-progress work, we treat these evaluation metrics as the primary protocol for the next phase rather than claiming final benchmark results.

\section{Experiments}

TODO


\subsection{Current Progress}

At the time of writing, we have completed the core environment formalization and implementation, with partially implementations of the stochastic kernels (proof and conjecture) and the query model prototype. We are currently integrating the environment with a multiagent training stack and validating kernel behavior via unit tests and small-scale rollouts.

\section{Discussion}
\subsection{Limitations}

We intentionally abstract away many aspects of real mathematical practice and formal theorem proving. The current environment does not generate or verify concrete Lean proofs; instead, it simulates proof and conjecture outcomes using stochastic kernels over a synthetic graph. As a result, empirical conclusions obtained in Forge may not transfer directly to a real theorem-proving market without additional validation.

A second limitation is the use of simplified truth labels, unique dependency realizations, and finite theorem universes. While these assumptions enable tractable experimentation and clearer credit assignment, they under-represent ambiguity, multiple proof strategies, and theorem statement expressivity in real formal systems.

Third, our initial query model is bucketed and low-capacity. Although this makes Bayesian updating simple and interpretable, it may fail to generalize well across heterogeneous theorem families or shifting prover capabilities. More expressive feature-based Bayesian models or learned uncertainty estimators are likely necessary for large-scale environments.

Finally, multiagent RL in partially observable, non-stationary market environments remains difficult. Even with CTDE, self-play, and curriculum learning, optimization may be unstable and outcomes may depend strongly on reward shaping and mechanism details. We therefore treat current training plans as exploratory rather than definitive.

\subsection{Next Steps}

Our immediate next steps are:

\begin{enumerate}
    \item \textit{Complete implementation and testing of environment kernels}, including property-based tests for monotonicity, admissibility gating, and consistency invariants.
    \item \textit{Integrate the environment with a CTDE training stack} and run small-scale experiments on procedurally generated instances.
    \item \textit{Implement market mechanism variants} (e.g., bounty-only, claims markets, hybrid designs) to study how incentives alter allocation efficiency.
\end{enumerate}

\section{Distribution of Work}
\subsection{Riyaz}

I will focus on the agent-learning interface and experimentation pipeline. This includes finalizing the structured trajectory tokenization scheme (observation/action/public-event tokens), implementing the prototype causal-transformer actor with pointer-style target selection heads, and integrating the centralized critic interface for CTDE training. I will also define the initial experiment matrix (curriculum stages, graph-generation parameter sweeps, and evaluation metrics) and run the first small-scale rollouts to validate policy behavior and logging quality.

\subsection{Shiv}

I will focus on environment robustness, testing, and training-stack integration. This includes completing and validating kernel implementations (unit and property-based tests for monotonicity, admissibility gating, consistency invariants, and dependency-closure publication), implementing the serialized observation/action/event API used by the transformer policy, and integrating the environment with the CTDE rollout runner. I will also implement and benchmark scripted baselines (feasibility-first proving, utility-greedy conjecturing, and publication heuristics) to generate offline warm-start trajectories for Stage 1 behavior cloning.


\bibliography{aaai23}
\end{document}
